% Papiers de Sophie Germain — Français 9117, lecture modernisée du volume entier.
%
% Pass: Opus 5 (claude-opus-5), 15 September 2026 — first pass, unchecked
% against the views by a human, and an interpretation rather than a record.
% Where this file and the transcriptions differ in substance, the
% transcriptions (batch-01.fr.tex … batch-06.fr.tex) are the record.
%
% Sources read: the six batch transcriptions of this volume and nothing else.
% No image was consulted, and no printed edition of the essay — neither the
% posthumous volume of 1833 nor any later reprint — nor any other
% transcription of it. Where the transcription has a gap, this reading says so
% and does not fill it. All six batches were transcribed on Opus 5; together
% they cover views 1–113, the whole volume.
%
% What the volume turned out to be. Not mathematics: one prose work, her essay
% « Considérations générales sur l'état des sciences et des lettres aux
% différentes époques de leur culture », in her hand from view 9 to view 110,
% a working copy corrected throughout and marked, in other hands, as copy for
% the press. Before it, the library's title leaf and the covering letter with
% which Lherbette gave the manuscript to the Bibliothèque royale. The spine is
% therefore single and continuous; the sections below follow her argument,
% and say where her own chapter heads (« Ch. I. », « Ch. 2 ») and the
% pencilled marks of other hands (« Ch 4 », « Ch. 5 », « Ch. 6 ») fall. No
% « Ch. 3 » stands on any leaf, and this reading does not place one.
%
% The prose standard. The body restates her argument in present-day French,
% in the third person, with short quotations of the leaf in her spelling.
% Each view's apparatus — folio, her struck page number, other hands,
% illegible and doubtful readings, deletions that show the thought changing,
% syntax she left unsettled — is carried in a footnote where that view's text
% begins in the restatement. Where she says something about the history of
% science or mathematics that is wrong or loose, the body says what is true
% and the footnote gives the leaf. Modern names (Laplace's determinism,
% Lessing's pregnant moment, the centre-of-mass theorem, the
% Lagrange–Dirichlet theorem…) are attached only where the statement
% coincides, and never as a claim that she read or cited them.
\documentclass[11pt,a4paper]{article}
% The preamble shared by every transcription of the mathematicians' manuscripts in Gallica.
%
% It rests on one idea: **uncertainty compounds, it does not resolve.** A
% transcription that smooths over a doubtful word has destroyed the only thing
% it was made for — a reader must be able to tell, at every point, what was on
% the page and what was supplied. The apparatus macros below are therefore the
% critical apparatus, not decoration.
%
% The same commands are recognised by `scripts/render.mjs`, which produces the
% site's reading view, and by `scripts/tei.mjs`. A macro added here must be
% added there too, or rendering fails loudly — which is the wanted behaviour.
%
% Comments here are English, like the rest of the repository. The documents
% this preamble sets are French, and that is deliberate: see the skills.

\ProvidesPackage{germain}[2026/09/15 Transcriptions of mathematicians' manuscripts in Gallica]

\usepackage{amsmath,amssymb,amsthm}
\usepackage{mathrsfs}
% Kept from the parent preamble so the renderer's diagram support stays
% available; these manuscripts draw no commutative diagrams, and nothing here uses it.
\usepackage{tikz-cd}
\usepackage[english,french]{babel}
\usepackage{fontspec}

% --- Greek ------------------------------------------------------------------
% The text face stays fontspec's default, Latin Modern, which has no Greek:
% XeTeX drops every glyph it lacks with a line in the log and nothing on the
% page, and the Greek words the manuscripts quote vanished from the PDFs. So
% the Greek and Greek Extended blocks get a XeTeX character class of their own,
% and entering or leaving a run of it switches the family to CMU Serif — the
% same Computer Modern design, with polytonic Greek — and back. Only the
% family changes: a Greek word in \textit or \textbf keeps its shape and series.
% The transitions to and from every other class carry the tokens that class
% has towards an ordinary letter, so French spacing before « ; » or inside
% « » still holds after a Greek word. No grouping, so no brace or \endgroup
% can come out unbalanced; maths is left alone, since XeTeX inserts nothing
% there. Kept identical in `exercices.sty`.
\makeatletter
\newfontfamily\germainGreekfont{cmun}[
  Extension=.otf, NFSSFamily=germaingreek,
  UprightFont=*rm, ItalicFont=*ti, SlantedFont=*sl,
  BoldFont=*bx, BoldItalicFont=*bi, BoldSlantedFont=*bl]
\newXeTeXintercharclass\germain@greek
\def\germain@greekrange#1#2{%
  \@tempcnta=#1\relax
  \loop
    \XeTeXcharclass\@tempcnta=\germain@greek
    \ifnum\@tempcnta<#2\advance\@tempcnta\@ne\repeat}
\germain@greekrange{"0370}{"03FF}
\germain@greekrange{"1F00}{"1FFF}
\def\germain@greekprev{\rmdefault}
\def\germain@greekin{%
  \edef\germain@greekprev{\f@family}\fontfamily{germaingreek}\selectfont}
\def\germain@greekout{\fontfamily{\germain@greekprev}\selectfont}
\def\germain@greekjoin{%
  \XeTeXinterchartokenstate=\@ne
  \@tempcnta=\z@
  \loop
    \ifnum\@tempcnta=\germain@greek\else
      \XeTeXinterchartoks\germain@greek\@tempcnta=\expandafter{%
        \expandafter\germain@greekout\the\XeTeXinterchartoks\z@\@tempcnta}%
      \XeTeXinterchartoks\@tempcnta\germain@greek=\expandafter{%
        \the\XeTeXinterchartoks\@tempcnta\z@\germain@greekin}%
    \fi
    \ifnum\@tempcnta<4095\advance\@tempcnta\@ne\repeat}
% babel-french sets its own classes, and saves and restores the interchar
% state, each time a language is selected: join again after it does.
\addto\extrasfrench{\germain@greekjoin}
\addto\extrasenglish{\germain@greekjoin}
\AtBeginDocument{\germain@greekjoin}
\makeatother

\usepackage{geometry}
\geometry{a4paper,margin=25mm,marginparwidth=28mm}
\usepackage{marginnote}
\usepackage{xcolor}
\usepackage[normalem]{ulem}
\setlength{\emergencystretch}{3em}
\usepackage{hyperref}
\hypersetup{colorlinks=true,linkcolor=black,urlcolor=black,pdfborder={0 0 0}}

% --- Batch metadata ---------------------------------------------------------
% Taken as they stand by the rendering script, which shows them at the head of
% the reading view. They fix what the file claims to cover. The macro names are
% the parent project's, so that the scripts stay shared; \folder here means the
% volume's slug — `fr-9115` — and \pages counts Gallica views.

\newcommand{\folder}[1]{\gdef\grothFolder{#1}}
\newcommand{\batch}[1]{\gdef\grothBatch{#1}}
\newcommand{\pages}[2]{\gdef\grothFirst{#1}\gdef\grothLast{#2}}
\newcommand{\foldertitle}[1]{\gdef\grothFoldertitle{#1}}

% The holder's own shelfmark, printed as they write it: « Français 9115 ».
\newcommand{\shelfmark}[1]{\gdef\grothShelfmark{#1}}

% The dating, copied verbatim from the holder's catalogue — never inferred
% here. « XIXe siècle » is the BnF's; a pass that dates a leaf from its content
% says so in a \note{}, as its own inference.
\newcommand{\dating}[1]{\gdef\grothDating{#1}}

% The status notice, printed in the title block and repeated as a diagonal
% watermark on every page of the PDF. The manuscripts are in the public
% domain; what this line declares is not a right but a quality — an
% unverified first pass by a machine. The skills write the line; a file
% without it gets no watermark, so leaving it out is visible in the output.
\newcommand{\watermark}[1]{\gdef\grothWatermark{#1}}

\gdef\grothFolder{?}\gdef\grothBatch{}\gdef\grothFirst{?}\gdef\grothLast{?}%
\gdef\grothFoldertitle{}\gdef\grothShelfmark{}\gdef\grothDating{}\gdef\grothWatermark{}

% --- The résumé -------------------------------------------------------------
\newenvironment{resume}
  {\small\setlength{\parskip}{0.45em}}
  {\par\medskip\hrule\medskip}

% --- Keywords ---------------------------------------------------------------
% In English, closing the modernised reading's résumé: the modern vocabulary
% under which to search for what the leaves construct. The manifest extracts
% them to tag the volume — this line is the single source of the tags.
\newcommand{\keywords}[1]{\par\smallskip\noindent
  {\footnotesize\textsc{Keywords} --- \textit{#1}}\par}

% --- The critical apparatus -------------------------------------------------

% The Gallica view number, in the margin. This is the anchor that lets a
% disagreement over a formula be settled by looking at *one* image rather than
% a volume of 750. The site uses it to turn the facsimile as the transcription
% is scrolled. A view is one scanned image: usually one side of a leaf.
\newcommand{\page}[1]{%
  \marginnote{\footnotesize\color{gray!70}\textsf{v.\,#1}}%
  \label{page:#1}\ignorespaces}

% The BnF's pencilled foliation, where the leaf carries it and a pass can read
% it: « 198r ». This is what the literature cites — « ff. 198r–208v » — and
% the one line that turns a citation into a view. Recorded, never inferred:
% a folio a pass cannot read is not written.
\newcommand{\folio}[1]{%
  \marginnote{\footnotesize\color{gray!50}\textsf{f.\,#1}}\ignorespaces}

% The modernised reading's coarser counterpart to \page{}: which views a
% section is drawn from.
\newcommand{\pagerange}[2]{%
  \marginnote{\footnotesize\color{gray!70}\textsf{v.\,#1--#2}}%
  \ignorespaces}

% Illegible: never guessed.
\newcommand{\ill}{\textcolor{red!60!black}{[\,\ldots\,]}}

% A doubtful reading: the word is given, and flagged as uncertain.
\newcommand{\uncertain}[1]{\uline{#1}}

% An editorial addition — a missing letter, an implied word.
\newcommand{\add}[1]{\textcolor{blue!50!black}{[#1]}}

% Struck out by the writer. What was crossed out is part of the document.
\newcommand{\struck}[1]{\sout{#1}}

% The transcriber's note, on the state of the page or the mathematics.
\newcommand{\note}[1]{\par\smallskip{\small\color{gray!60!black}\textit{[#1]}}\par\smallskip}

% A marginal note of hers, distinct from the body of the text.
\newcommand{\marginal}[1]{\par\smallskip\noindent\hspace{1em}%
  {\small\color{gray!60!black}#1}\par\smallskip}

% --- Title ------------------------------------------------------------------
% The title block is French, like the documents themselves.

\AddToHook{shipout/background}{%
  \ifx\grothWatermark\empty\else
    \begin{tikzpicture}[remember picture, overlay]
      \node[rotate=52, scale=3.4, align=center, text=gray!18,
            font=\sffamily\bfseries]
        at (current page.center) {\grothWatermark};
    \end{tikzpicture}%
  \fi}

\AtBeginDocument{%
  \begin{center}
    {\large\bfseries Manuscrits de mathématiciens (Gallica) ---
      \ifx\grothShelfmark\empty\grothFolder\else\grothShelfmark\fi}\\[2pt]
    {\itshape\grothFoldertitle}\\[4pt]
    {\small vues \grothFirst--\grothLast
      \ifx\grothBatch\empty\else\ (lot \grothBatch)\fi}\\[2pt]
    \ifx\grothDating\empty\else
      {\small\color{gray!60!black}Datation du catalogue : \grothDating}\\[2pt]
    \fi
    \ifx\grothWatermark\empty\else
      {\small\bfseries\color{red!45!gray}\grothWatermark}\\[2pt]
    \fi
    {\footnotesize\color{gray!60!black}Transcription établie d'après les images IIIF de Gallica.
     Source gallica.bnf.fr / Bibliothèque nationale de France.\\
     Les lectures incertaines sont soulignées, les illisibles notées [\,\ldots\,],
     les ajouts entre crochets ; \textsf{v.} numérote les vues de Gallica,
     \textsf{f.} la foliotation au crayon de la BnF.}
  \end{center}
  \vspace{1em}}

\endinput


\folder{fr-9117}
\shelfmark{Français 9117}
\pages{1}{113}
\dating{1801-1900}
\watermark{Lecture automatique\\interprétation non vérifiée}
\foldertitle{« Considérations générales sur l'état des sciences et des lettres aux différentes époques de leur culture » par Sophie GERMAIN — lecture modernisée du volume entier}

\begin{document}

\section*{Résumé}

\begin{resume}
Ce volume ne contient pas une ligne de calcul. C'est le manuscrit de travail,
de sa main, d'un essai de philosophie : les « Considérations générales sur l'état
des sciences et des lettres aux différentes époques de leur culture ». Une
centaine de pages, biffées et récrites entre les lignes, portant en marge des
noms d'autres mains qui semblent ceux des compositeurs d'une imprimerie ; en tête, le feuillet de la Bibliothèque et la
lettre par laquelle son neveu Lherbette donna le manuscrit, après l'impression
posthume du texte. Sophie Germain y écrit en géomètre sur tout ce qui n'est pas
la géométrie.

Le problème qu'elle pose est ancien, et elle le pose de façon précise. Nous
jugeons partout — du vrai dans les sciences, du beau dans les arts, du bien en
morale — d'après un même sentiment d'unité, d'ordre et de proportion. Ce
« type », qu'elle croit inné, dit-il quelque chose de l'être lui-même, ou
seulement de la manière dont un esprit humain peut concevoir les choses ? Elle y voit
comprise la question de Kant : notre logique est-elle celle de la raison
absolue, ou seulement de la raison humaine ?

Son idée tient en trois temps. D'abord, le poète et le géomètre font la même
chose : un tâtonnement, une idée simple, un plan, une exécution, des corrections
où le tact du géomètre qui choisit les formules à écrire n'est autre chose que le
goût. Ensuite, l'histoire de l'esprit — des dieux qu'on imagine à son image
jusqu'à Descartes et Newton — est celle d'un type du vrai qui s'est toujours fait
sentir mais qui, étant abstrait, ne pouvait dire quelles réalités particulières
existent : d'où l'esprit de système, l'alchimie, les tourbillons, puis
l'observation, le « comment », le « combien », et la physique devenue science
exacte. Enfin, la réponse : la langue du calcul, qui tire d'une seule réalité
toutes celles qui lui sont liées, et que l'expérience vient confirmer par une
autre voie, prouve que l'ordre cherché appartient à l'être ; le monde, mieux
connu, paraîtrait nécessaire, et le contingent n'est que ce dont nous ignorons la
nature. Elle en tire une confiance dans l'analogie, qu'elle pousse loin : la
mécanique — équilibre stable et instable, impulsion hors du centre de gravité,
chocs, levier — lui sert à décrire les sociétés et les révolutions ; et la
musique, « toute métaphysique », lui apparaît comme une langue qui a « ses
phrases, ses périodes, ses repos, ses hardiesses ».

Ce qu'elle voulait faire, elle l'a écrit à l'endroit où l'essai tourne : « Les
préliminaires qu'on vient de lire nous ont paru nécessaires pour bien entendre
les idées que nous allons exposer, ils en fixeront le sens et serviront peut être
à leur faire pardonner ce qu'elles semblent avoir de hardiesse et de nouveauté. »
Et elle finit sur une promesse : que les poètes montreront un jour l'univers « tel
qu'il est réellement », et que le génie verra que l'art de créer n'était que celui
de copier les parties d'un tableau qu'il pourra enfin peindre en entier.

Plusieurs de ses affirmations de mécanique et d'histoire des sciences doivent être
redressées, et la lecture le fait à chaque endroit : une impulsion donnée hors du
centre de gravité ne perd rien pour le mouvement d'ensemble ; les perturbations
ne s'éteignent que dans les systèmes dissipatifs ; Archimède fondait la statique,
non la dynamique ; Bacon précède Descartes ; Fermat a eu part à la géométrie
analytique. Les noms modernes sous lesquels chercher : le déterminisme de Laplace
et la probabilité comme mesure de l'ignorance ; la philosophie critique de Kant et
le statut des formes \emph{a priori} ; l'efficacité des mathématiques dans les
sciences de la nature ; le \emph{Discours préliminaire} de d'Alembert et l'unité
du savoir ; la stabilité de l'équilibre (Lagrange–Dirichlet) et le théorème du
centre d'inertie ; le « moment fécond » de Lessing, et, pour la musique, le débat
sur l'expression des sentiments.

\keywords{philosophy of science, unity of knowledge, Kant's critique of reason,
a priori knowledge, necessity and contingency, Laplacian determinism,
probability as ignorance, analogy in scientific discovery, history of
mathematics, analytic geometry, Newtonian mechanics, stability of equilibrium,
centre of mass theorem, mechanical metaphors in political thought, aesthetics,
music as language, d'Alembert's Preliminary Discourse}
\end{resume}

\section*{Ce que le volume contient}

\pagerange{1}{113}
Cent treize vues, cinquante-quatre feuillets numérotés par la Bibliothèque.
Les quatre premières vues sont la mire du microfilm, la reliure et une garde ;
les trois dernières, une page blanche et le plat inférieur.\footnote{Vue 1,
mire « Ce document a été microfilmé tel qu'il a été relié » ; vues 2 à 4, plats
marbrés à l'étiquette « FR. 9117 », contre-plat, garde blanche ; vue 6,
blanche ; vue 111, feuillet blanc ; vues 112 et 113, reliure. Rien de cela n'est
transcrit.} Entre les deux, trois pièces.

\begin{itemize}
  \item \emph{Le feuillet de titre de la Bibliothèque} (f. 1, vue 5), d'une main
    de bibliothécaire : l'ancienne cote « Supplt fr. 2209 », la mention
    « Manuscrit autographe de Mademoiselle Germain », le titre de l'édition de
    1833, « donné par M. Lherbette député \& neveu de l'auteur », et une
    référence d'entrée « R.B. N° 106 4. 1833 ».
  \item \emph{La lettre d'envoi} de Lherbette au conservateur Champollion
    (f. 2, vues 7 et 8), qui n'est pas de sa main.
  \item \emph{L'essai lui-même}, de la vue 9 à la vue 110, de sa main : quatre-vingt-dix-neuf
    pages écrites, sur les rectos et les versos, trois versos seulement étant
    restés blancs.\footnote{Versos blancs : vue 12 (dos du feuillet collé de
    la vue 11), vue 14 (dos de la vue 13) et vue 66. Le feuillet d'addition de
    la vue 11 est plus petit que les autres et monté.}
\end{itemize}

L'essai n'est pas une copie au net. C'est un manuscrit de travail : chaque
page est biffée et récrite entre les lignes, des alinéas entiers sont barrés
d'une croix et repris dessous, deux passages portent en marge « Variante »,
et bien des alinéas s'ouvrent par un crochet de sa plume. Les corrections
s'éclaircissent vers la vue 33 et redeviennent serrées dans les pages sur la
musique. Deux fois seulement elle écrit elle-même un titre de chapitre :
« Ch. I. » à la vue 9, « Ch. 2 » à la vue 23 — et le titre de ce second
chapitre est celui que la Bibliothèque a donné au volume.

Deux numérotations courent sur les feuillets. En haut à droite des rectos,
d'un trait fin et gris, la foliotation de la Bibliothèque, de 1 (vue 5) à 54
(vue 109) ; elle est citée ici « f. », recto seulement, et n'est jamais
suppléée.\footnote{Le folio 3 manque à la série, qui passe de 2
(vue 7) à 4 (vue 9) ; la transcription du lot 1 suppose qu'il répondrait au
panneau d'adresse de la vue 8, qui ne porte aucun chiffre visible ; à la vue 71 le second chiffre est ambigu, et la
transcription n'y porte aucun folio. Toute la série a été lue sur les vues,
aucune n'est inférée des numéros de vue.} En haut à gauche de chaque page, à
l'encre et presque toujours barrée d'un trait, sa propre pagination, de 1 à
100 ; elle est donnée en note, page par page, là où la transcription l'a
lue.\footnote{Aucun « 4 » n'a été trouvé dans cette série ; le feuillet collé
de la vue 11 porte seulement une croix. Dans le lot 4 (vues 61 à 80) les
nombres sont presque tous illisibles sous la biffure : on lit 60, 66 et 69 aux
vues 70, 76 et 79.}

D'autres mains ont passé sur le manuscrit. Des noms soulignés, d'une écriture
plus grande ou plus petite, se lisent en marge ou en tête des alinéas :
« Gutmann » (vues 9 et 69), « Cherenoz » (vue 17), « Deraihe », « Deracke »,
« Deraché », « Derache » (vues 19, 55, 73, 106), « Pennequin » ou « Pemeguin »
(vues 39, 79, 101), « Varlé » (vue 61), « Chévenon » (vue 92), « Georges »
(vue 99), la plupart lus avec doute. Au crayon, d'une autre main, « Ch 4 » en
tête de la vue 60, « Ch. 5 » à la vue 79, « Ch. 6 » et « Avenir de » à la
vue 84 ; à l'encre noire épaisse, un « II » romain sous le folio 16 (vue 33).
La transcription y voit, et le dit inférence, les noms des compositeurs
marquant où chacun reprenait la copie, et les divisions d'une mise sous
presse.\footnote{Autres marques de même genre : « N. 2 » encadré (vue 21),
« Fl. 7 » encadré (vue 69), « Double » au crayon (vue 89), un « avait » en
lettres droites dans l'interligne (vue 70), et, peut-être, des additions d'une
écriture plus fine à la vue 30. Leur rôle n'est établi nulle part dans les
transcriptions, et cette lecture ne l'établit pas davantage : elle n'a
confronté le manuscrit à aucune édition imprimée.} Il n'y a de « Ch. 3 » sur
aucun feuillet. Cette lecture ne le place donc pas ; ses sections suivent
l'argument, et disent au passage où tombent les marques.

Une seule forme grammaticale du texte dit qui écrit, et c'est un accord :
« La digression historique à laquelle nous nous sommes livrée ».\footnote{Vue
51. Le participe est au féminin sur le feuillet.} Aucune date n'est écrite
dans l'essai. Ce qu'il nomme de plus récent — Kant « de nos jours », les
« successeurs de Newton », l'entrée des éléments de l'algèbre dans
l'éducation — ne permet pas de le dater mieux que la notice ne le fait, et
cette lecture ne s'y essaie pas.

\section*{La lettre d'envoi}

\pagerange{5}{8}
Un petit feuillet monté sur onglet, d'une main qui n'est pas celle de Sophie
Germain.\footnote{Vue 7, f. 2r ; timbre « Bibliothèque royale » non transcrit.
L'adresse est au verso, vue 8 : « Monsieur Champollion Conservateur à la
Bibliothèque Royale ».} Son auteur, « D'après le désir que vous avez bien voulu
m'en témoigner », adresse au conservateur, avec un exemplaire d'un « ouvrage
posthume de Mlle Germain », le manuscrit, à joindre s'il le juge à propos à
ceux que la Bibliothèque conserve déjà. En post-scriptum, il demande la
permission de joindre un second exemplaire pour Letronne, qui, ne connaissant
pas Mlle Germain, ne comprendrait pas un envoi direct.\footnote{« ne connaissant
pas \emph{personnellement} Mlle Germain » et « l'envoi que je lui \emph{en}
ferais » : les deux mots en italique sont des lectures incertaines de la
transcription.}

La lettre est datée « 7bre. », sans jour ni année, et signée
« Lherbette ».\footnote{La signature est lue d'après la notice du feuillet de
titre (vue 5), qui nomme le donateur « M. Lherbette député \& neveu de
l'auteur ». La date est écrite à gauche de la formule finale.} Deux
identifications sont de cette lecture, non du feuillet. Le conservateur
Champollion de la Bibliothèque royale est vraisemblablement Jacques-Joseph
Champollion-Figeac, conservateur des manuscrits depuis 1828 — son frère
Jean-François, le déchiffreur, n'était pas de la Bibliothèque et mourut en mars
1832. L'« ouvrage posthume » est l'essai lui-même, imprimé après la mort de
Sophie Germain (27 juin 1831) ; la notice du feuillet de titre en donne le
titre et porte « 1833 » dans sa référence d'entrée, ce qui situe vraisemblablement
l'envoi en septembre 1833, sans que rien ne l'écrive.\footnote{La lecture
« septembre 1833 » combine le « 7bre » de la lettre et le millésime de la
référence d'entrée ; c'est une inférence. Letronne n'est pas autrement nommé :
Jean-Antoine Letronne, l'helléniste, est le seul de ce nom que le contexte
suggère, et la lettre ne le qualifie pas.}

\section*{Chapitre I. « Comment les sciences et les lettres sont dominées par un sentiment qui leur est commun »}

\pagerange{9}{22}
Le premier chapitre soutient une seule thèse, et la prouve par une
comparaison. La thèse : toutes les œuvres de l'esprit, science ou poème,
obéissent à un même sentiment d'ordre, de proportion et de simplicité, qui est
pour nous la marque du vrai. La comparaison : suivre pas à pas le poète et le
géomètre, du premier essai à la dernière correction, et constater qu'ils font
la même chose.\footnote{Vue 9, f. 4r ; page 1 de sa numérotation. Au centre,
« Ch. I. », souligné, puis le titre. En marge du premier alinéa, d'une autre
écriture, « Gutmann », souligné ; l'alinéa est ouvert par un crochet.}

\subsection*{Le type : ordre, proportion, simplicité}

\pagerange{9}{13}
À considérer ensemble les travaux de l'esprit humain, on est frappé de leur
ressemblance. Partout certaines lois ont été suivies ; là où elles ne l'ont
pas été, le défaut s'en est fait sentir, et l'ouvrage — traité de doctrine ou
simple divertissement — n'a pas duré : la curiosité épuisée, l'oubli. Ces lois
ont gouverné la pensée bien avant la réflexion. Le spectacle de l'univers en
portait l'empreinte, la mémoire les a reproduites, l'imagination leur est
restée assujettie jusque dans ses caprices, et la raison, plus tard, en a fait
son guide.\footnote{« soumise », biffé, remplacé par « assujetie ».}

Faut-il y voir les conditions mêmes de l'être, ou seulement la manière dont
nous pouvons concevoir le possible ? Elle pose les deux branches de
l'alternative — le « type » que nous trouvons en nous et dans les choses
« nous révèle les conditions de l'être », ou bien, nous appartenant en propre,
« atteste seulement la manière dont nous pouvons comprendre les possibles » —
et déclare la question hors de notre portée. Il suffit à son propos de
constater comment un sentiment profond d'ordre et de proportions devient pour
nous, en toute matière, le caractère du vrai, et qu'ainsi, dans les genres
d'études les plus divers, nos recherches tendent au même but par les mêmes
procédés.\footnote{Fin de la vue 9 et haut de la vue 10 (page 2) : « pouvons »,
biffé, remplacé par « pourrons », lecture incertaine ; le premier mot biffé de
la vue 10 est noirci et illisible ; « comment », biffé, est remplacé par
« que ». La question qu'elle écarte ici est celle qu'elle reprend, en la
rapportant à Kant, à partir de la vue 35.}

Le plan d'un ouvrage, l'argument d'un poème, exigent de la clarté ; les
parties doivent se lier assez bien pour qu'on en saisisse le rapport d'un coup
d'œil ; l'esprit veut un ordre facile et se plaît à la simplicité, « source de
l'élégance en tout genre ». Le merveilleux lui-même y est soumis : là où la
fiction remplace le réel, un certain « module intellectuel » tient lieu de ce
qui manque à la réalité des objets. Les oracles du goût et les arrêts de la
raison se ressemblent, parce qu'un même type universel appartient au beau et au
vrai. L'histoire naturelle le montre autrement : nous classons les êtres selon
nos convenances, et la notion méthodique des genres et des espèces imprime à
cette science le cachet de l'esprit humain.\footnote{« d'un coup
d'\emph{œil} » : lecture incertaine. Le mot « module », qui reviendra dans les
pages sur les arts (vues 89 à 92, 101, 110), a ici son sens d'architecte :
l'unité de mesure propre à un ordre, à laquelle toutes les proportions se
rapportent. C'est une interprétation de cette lecture ; la transcription
laisse le sens du mot en suspens.}

Dans les sciences exactes, le sentiment d'ordre et de proportion cède la place
à la connaissance certaine d'un ordre déterminé et de proportions mesurables.
L'esprit semble y avoir changé de marche ; en réalité, la ressemblance à son
modèle intérieur n'y est plus seulement pour lui le caractère du vrai : il
l'atteint de plus près, l'objet réalise au plus haut degré ce qu'il cherche
partout ailleurs, et son attention s'y absorbe tout entière.\footnote{La phrase
commence au bas de la vue 10 et s'achève en tête de la vue 13 (f. 6r, page 3) ;
la vue 11 est un feuillet d'addition, et la vue 12 son verso blanc.}

Sans doute la lecture d'un roman ne ressemble pas à l'étude d'un traité de
géométrie, et certains esprits, livrés tout entiers aux images ou tout entiers à
la démonstration, deviennent incapables de l'autre genre. Il ne faut pas en
conclure qu'aucun lien n'unit ces œuvres. Qu'on assiste à leur création : on
verra l'esprit guidé partout par la prévision de certains résultats, suivant
une méthode constante, et il deviendra évident que la littérature et les
découvertes des sciences ont été inspirées par un même sentiment d'ordre et de
proportion, « régulateur de tout mouvement intellectuel ».\footnote{Vue 13. La
fin de cette phrase n'est pas établie au manuscrit : après « la littérature la plus »
vient un mot illisible ; après « et les découvertes », un mot recouvert
commençant par « d », deux mots biffés lus avec doute « ce genre »,
« s'enrichit » écrit au-dessus, puis un mot surchargé devant « sciences ». La phrase
modernisée ne garde que ce qui en est sûr ; l'épithète de la littérature est
perdue.}

Ce qui plaît dans un trait de génie le confirme. Dans les sciences, dans les
beaux-arts, dans les lettres, il plaît pour une seule et même raison : il
découvre d'un coup une foule de rapports qu'on n'avait pas aperçus. On se
trouve transporté dans une région élevée d'où l'on voit un ordre inattendu
d'idées ou de sentiments ; la surprise émeut l'âme, qui rend un hommage
involontaire à son bienfaiteur, et cet hommage est pour elle un plaisir de plus.
L'esprit humain obéit à des lois qui sont celles de sa propre existence ; elles
lui donnent une mesure commune entre toutes les existences qu'il conçoit, et
deviennent le mobile de tous ses travaux et la source de tous ses
plaisirs.\footnote{Ce paragraphe est le feuillet collé de la vue 11 (f. 5r),
plus petit, marqué d'une croix « + » ; la même croix est tracée au bas de la
vue 13, sous le dernier alinéa, et c'est là que cette lecture le place, alors
que la transcription le laisse dans l'ordre des vues. L'ordre des deux phrases
est inversé ici : le feuillet commence par les lois de l'esprit et finit sur le
trait de génie. Corrections : « l'éloquence », biffé, devient « littérature » ;
« conduits » devient « transportés » ; « nouveau » devient « inattendu ».}

\subsection*{Le poète et le géomètre, du premier essai à la dernière correction}

\pagerange{15}{19}
Premier temps, les essais. Le sujet choisi, les idées se pressent chez le
poète ; mille ressorts semblent pouvoir animer la composition ; il en suit un,
y renonce, en choisit un autre, complique son mécanisme, n'en est pas content,
revient sur ses pas — et du milieu de ce tumulte surgit une idée simple, dont
il sent que c'était elle qu'il cherchait. Le géomètre, lui, part d'une
remarque ou d'un fait inattendu ; il examine ce que la science faite peut lui
fournir et circonscrit son sujet ; il entrevoit des résultats qu'il ne peut
encore atteindre, s'élance, craint de s'être égaré, doute, rétrograde ; les
idées adventices compliquent le sujet et suspendent le jugement ; mais à
travers ce chaos il distingue une idée simple, et sait qu'elle sera
féconde.\footnote{Vue 15, f. 7r ; page 5 de sa numérotation. Un mot biffé
illisible après « il ». « l'avoit d'abord guidé » est ainsi au manuscrit.}

Deuxième temps, le plan. Le poète ne perdra pas de vue l'idée principale ;
elle donnera à l'ouvrage l'unité d'intérêt et d'action, « source de toute
beauté véritable », et lui permettra de satisfaire à ce besoin d'ordre et de
proportion qu'un sentiment universel a mis au premier rang des préceptes du
goût et de la raison. Le géomètre, de son côté, emploie toutes ses forces à
dérouler la chaîne des vérités contenues dans sa vérité première, et nulle part
l'unité de composition n'est plus sensible : l'ordre de son travail est
déterminé, il ne pourrait l'intervertir ; l'évidence est la condition de son
succès ; il choisit la méthode qu'il croit propre à l'y conduire, et entre
avec joie dans la carrière ouverte à ses espérances.\footnote{Vue 16, page 6.
Le premier alinéa de la page est barré d'une grande croix, biffé ligne à ligne,
et récrit aussitôt dessous ; il contient quatre mots illisibles et ne diffère
de sa reprise que par l'ordre des idées. Dans la reprise, les corrections se
chevauchent et la phrase reste incomplète au manuscrit : entre « elle lui
offrira le moyen » et « besoin d'ordre et de proportion » le verbe manque ;
« satisfaire à ce besoin » est supplé ici par le sens. Mots incertains :
« dérouler » (la chaîne des vérités), « au », et, biffés, « amenaient » et
« avant » (en marge). Vue 17, f. 8r, page 7 : avant « Et il entre alors avec
joie », deux lignes biffées disaient que « l'autorité de sa démonstration
succèdera aux simples prévisions de son jugement » — la phrase sur la
démonstration a été retirée, et la joie de l'entrée en carrière est restée.}

Troisième temps, l'exécution. Les divisions adoptées gardent entre elles une
juste proportion d'étendue, qui ménage à l'attention des repos sans rompre la
continuité. Pour remplir les cadres, les deux auteurs s'abandonnent de nouveau
à leur génie ; mais, les limites du sujet étant fixées, ils ne craignent plus
de s'égarer, l'un dans le champ d'une imagination inventive, l'autre dans
« cet océan des possibilités d'où on aborde avec tant de difficulté sur le
terrain ferme de la vérité démontrée ». Il leur viendra encore des idées nées
du sujet qui nuiraient à la rapidité ou à la clarté ; ils ne les écartent pas
d'abord, de peur d'arrêter l'élan. Plus tard, relisant leurs ébauches, ils
n'en gardent que les traits nécessaires, et deviennent les juges de leur propre
ouvrage.\footnote{Vue 17. En marge de l'alinéa « Les auteurs dont nous
comparons », d'une autre écriture, un nom lu avec doute « Cherenoz ». Le début
de la phrase « Il se présentera toujours encore des idées » est une lecture
incertaine d'une ligne surchargée, précédée d'une ligne biffée (« Comment
l'époque ou et époque des travaux ») et de deux groupes de mots biffés
illisibles ; « dans le cours du travail », entouré d'un trait, n'a pas de place
sûre.}

Quatrième temps, les corrections. On examine d'abord la marche des idées :
celles qui partageraient l'intérêt ou suspendraient l'attention sont déplacées,
et vont enrichir un épisode ou une note savante, ou deviennent l'origine d'un
autre ouvrage — « ainsi la branche développée dans la saison actuelle offre
quelquefois le rudiment d'une végétation prochaine ». Vient ensuite le style.
L'homme de lettres pèse le choix et l'ordre des mots, l'harmonie du vers ou de
la phrase, sous le jugement d'un goût tantôt prompt, tantôt lent, qui semble
n'obéir qu'à ses caprices et suit pourtant toujours les règles de la
raison.\footnote{Vue 18 ; en tête, un chiffre sous une tache, sans doute 8. Le
mot « style » est en partie couvert d'une tache et lu avec doute.}

Or « la langue des calculs » a aussi son style, et tous les auteurs ne
l'écrivent pas également bien. Au choix des mots correspond celui des
caractères : ceux-ci sont conventionnels au point qu'il faut chaque fois dire
quelle valeur on leur donne, et leur emploi obéit pourtant à des convenances
qui ne tiennent pas seulement à l'habitude. Les formules remplacent la phrase
et peuvent être plus ou moins élégantes. L'analyse parle aux yeux : là où la
musique demande l'accord des sons, elle doit présenter entre ses éléments des
rapports d'ordre et de simplicité saisissables d'un coup d'œil. Les initiés
trouvent dans la contemplation des formules un charme qui les attire à l'étude ;
et le tact qui fait choisir aux bons auteurs quelles formules écrire et
lesquelles seulement indiquer, avec des décisions tantôt rapides, tantôt
réfléchies, n'est autre chose que le goût, appliqué à des objets qu'on croyait
étrangers à son empire.\footnote{Vue 18 (fin) et vue 19, f. 9r, page 9. En tête
de la vue 19, d'une autre écriture, un nom lu avec doute « Deraihe ». Les
corrections de « Les personnes \ldots{} de charme » se superposent en trois
couches ; une première rédaction, biffée, disait que les initiés « trouvent
dans l'analyse les », un mot illisible, « une sorte d'attrait », et un second mot illisible.
« leur », devant « fait choisir », est couvert d'une tache. Le mot
« harmonie », sur le feuillet, désigne l'accord entre les sons ; « musique » est
de cette lecture.}

Le résumé qu'elle fait du chapitre tient en une phrase : les productions de
l'esprit les plus diverses ont entre elles de vraies ressemblances, et le même
sentiment d'ordre et de proportion préside à l'inspiration, guide son emploi,
et se fait sentir jusque dans les dernières corrections de l'ouvrage
achevé.\footnote{Vue 19. La ligne qui suit, biffée, s'interrompt au bas du
feuillet (« tellement que dans des genres de perceptions »).}

\subsection*{Pourquoi on les a crus séparés}

\pagerange{20}{22}
Si la marche de l'esprit est partout la même, les objets qu'il peut envisager
sont d'une variété infinie, et au premier coup d'œil c'est cette variété qui
frappe, plus que l'identité des rapports. Aussi des opérations
intellectuelles identiques ont-elles reçu des noms différents selon leurs
sujets. Et la différence des mots — d'autant plus naturelle que chaque branche
du savoir a longtemps été exclusive de toutes les autres — perpétue l'idée
d'une séparation réelle entre les facultés de l'esprit : comme si l'allégorie
n'était pas soumise aux préceptes de la raison, et comme si la découverte
d'une loi de la nature avait pu se passer de l'imagination.\footnote{Vue 20,
page 10 ; aucun folio (verso). Le haut de la page est très raturé, et la phrase
« Au premier coup d'œil \ldots{} l'identité des rapports » est reconstituée par
la transcription à partir de corrections superposées dans deux interlignes ;
l'ordre retenu est une conjecture, un mot surchargé n'est pas lu, et « nous
avons » (dans « cette variété dont nous avons parlé ») n'est pas écrit : c'est
un ajout de l'éditeur. Mots incertains : « confondues » (les opérations
« confondues, les mêmes »), « de l'esprit ».}

C'est que les auteurs effacent la trace de leur travail. Le poète ne dira pas
les discussions subtiles qui ont précédé le choix de ses emblèmes ; l'homme de
génie qui a surpris un secret de la nature ne dira pas combien son imagination
s'est égarée autour de la route qui devait le conduire à la connaissance
certaine d'une vérité qu'il est à présent en état de démontrer. Chacun fait
disparaître ses premiers essais pour ne garder que les formes propres au sujet.
Le lecteur, lui, cherche selon son humeur un délassement ou une instruction ;
le titre du livre l'avertit du degré d'attention qu'il aura à fournir ; il est
donc porté à croire que les auteurs ont écrit, les uns dans l'abandon d'une
imagination qui erre en liberté, les autres avec la méthode austère d'une
déduction sans écart. De là cette séparation, jadis si respectée, entre le
domaine de l'imagination et celui de la raison.\footnote{Fin de la vue 20
(« l'adoption » des emblèmes, et « a » devant « choisis », lectures incertaines)
et vue 21, f. 10r, page 11 barrée. Le haut de la vue 21 est très raturé, et la
syntaxe du passage n'est pas arrêtée au manuscrit : « Le lecteur vient ensuite
donc cherche » ; « de n'avoir à faire \emph{usage} » est incertain ; quatre
mots biffés sont illisibles. En marge, « N. 2 » dans un cadre. Le dernier
alinéa de la vue 21, « Le bienfait de l'imprimerie », est barré d'une grande
croix et récrit au verso.}

Une autre cause a joué. Dans des temps éloignés, l'extrême division du travail,
nécessaire à la science naissante, avait accrédité l'idée que les facultés de
l'âme sont spécialisées. Aujourd'hui l'imprimerie assure à l'esprit humain la
jouissance de tout ce que les générations précédentes ont accumulé
d'observations, de comparaisons, de théories et de vérités incontestables ; il
n'a plus à refaire les premiers pas, ses forces augmentent chaque jour, et
l'instruction approfondie le ramène par une voie sûre vers ces idées de
simplicité et d'unité qui furent autrefois des révélations du génie, devinant
sa propre nature et cherchant à en étendre les lois à l'univers.

Conclusion du chapitre : les sciences, les lettres et les beaux-arts ont été
inspirés par un seul et même sentiment ; chacun, avec ses moyens propres, a
produit des copies sans cesse renouvelées d'un même modèle inné, « type
universel de vérité ». Le chapitre suivant, annonce-t-elle, montrera par
l'histoire de l'esprit humain comment, jusque dans ses écarts, tous ses efforts
ont tendu vers l'ordre, la simplicité et l'unité.\footnote{Vue 22, page 12
barrée (verso du f. 10). Deux mots chargés d'encre, au-dessus de « augmentent »
et devant « révélations », ne sont pas lus. La construction « aujourd'hui que le
bienfait de l'imprimerie, en assurant \ldots{}, il n'aura plus » est restée
telle au manuscrit ; la phrase modernisée la dénoue. « l'ame » remplace
« l'esprit », biffé.}

\section*{Chapitre 2. « Considérations générales sur l'état des sciences et des lettres aux différentes époques de leur culture » : les premiers âges de la pensée}

\pagerange{23}{33}
Le second chapitre porte le titre de l'essai entier. Il annonce une histoire,
et prévient aussitôt qu'il ne s'astreindra à « aucun ordre historique » : c'est
une genèse conjecturale, qui reconstruit la marche la plus simple que l'esprit
ait pu suivre, non une chronique.\footnote{Vue 23, f. 11r, page 13 barrée.
« Ch. 2 » est écrit au-dessus du titre et souligné. En marge du premier
alinéa, « Variante », suivi d'un mot biffé : le chapitre s'ouvre donc sur une
rédaction que l'auteur tenait pour une variante. La formule « Sans nous
astreindre à aucun ordre historique » est à la vue 26.}

\subsection*{L'origine de la littérature, et l'homme qui se prend pour modèle}

\pagerange{23}{26}
La littérature a commencé le jour où l'homme, sortant du cercle de ses
intérêts, a voulu communiquer à ses semblables des sentiments et des idées sans
but utile. Le récit des événements remarquables et la peinture des grandes
scènes de la nature n'étaient encore que des copies. Lorsque l'homme de génie
parvint à reproduire des impressions à l'aide d'une action dont il avait
imaginé les ressorts, il s'était déjà élevé à la notion abstraite de l'ordre,
pour en tirer la première règle de la composition : voulant disposer de
l'attention des autres, il pratiqua l'unité d'action, l'unité d'intérêt et la
clarté d'exposition comme des moyens de succès, bien avant que l'esprit
d'examen en fît des préceptes de l'art.\footnote{Vue 23. « nous trouvons
qu'elle » (la littérature a commencé), écrit au-dessus de la ligne, est une
lecture incertaine, et la phrase n'a pas été remise d'accord avec cet ajout ;
« de » est porté en marge devant « La littérature », dont la majuscule reste.
Cinq mots biffés dans ce paragraphe sont illisibles, et une addition de six
lignes, dans la marge de gauche, est entièrement biffée et ne se lit pas. Le
bas de la page est laissé blanc.}

Jeté au milieu d'une infinité de merveilles, l'homme n'a rien trouvé de plus
merveilleux que lui-même. Il a étendu son existence à tout ce qui l'entourait :
connaissant d'abord son individualité, cherchant partout sa propre image, il a
personnifié les êtres inanimés et fait présider à tous les phénomènes des êtres
intellectuels nés de son imagination. Se manifestaient déjà, à ce premier âge,
le sentiment profond d'un lien commun entre tous les êtres et celui d'un type
universel, empreint dans l'intelligence pour lui servir de modèle. Les sciences
n'existaient pas, mais le besoin d'expliquer se faisait sentir : la première
littérature fut poétique, et ce qui tenait lieu de physique ne l'était pas
moins ; les deux branches, si séparées aujourd'hui, étaient alors
confondues.\footnote{Vue 24, page 14 barrée (verso du f. 11). Un mot biffé
devant « individualité » est illisible ; « ont » (présidé) est incertain. De
petites croix de renvoi suivent « imagination », « sentir », « poétique » et
« premiers tems », et un « 1 » est porté au-dessus de « êtres inanimés », sans
objet sur la page ; la même chose se répète aux vues 25, 31, 32, 37, 38, 39 et
40.}

Qu'importait, pour la composition, que le sujet fût l'homme ou un dieu, un
demi-dieu, un génie doté des passions humaines ? Des êtres si pareils
pouvaient agir ensemble, le merveilleux les unissait. On voit dans ces premiers
essais le goût des idées générales et le sentiment d'analogie, qui
reparaîtront sous mille formes. L'homme se connaissait intelligent, dirigeant
ses actions vers un but ; il voit dans la nature un ordre et une succession qui
semblent tendre vers un but ; il n'imagine d'autre cause qu'une intelligence et
une volonté, et ne peut les concevoir sans un être qui les porte. Ces êtres,
il ne les voit pas : il les suppose invisibles — dieux, demi-dieux ou génies
subalternes selon l'importance de leurs actes —, amis ou ennemis, faits à notre
image ; et comme on ne peut ni les voir, ni les entendre, ni les toucher, ils
sont immatériels : ce sont des esprits. En marge, elle résume : « Dès qu'il
porte ses regards autour de lui que cherche-t-il ? ce qu'il a trouvé en lui
même. »\footnote{Vue 25, f. 12r, page 15 barrée. L'addition marginale est
appelée par une croix après « sa propre existence » ; son « que » est
incertain. « d'autre » (il n'imagine d'autre cause) est incertain, sa première
lettre empâtée ; au-dessus de « n'imagine », un mot barré, peut-être
« suppose ».}

Fidèle à sa pensée constante, l'homme a toujours tenu sa propre existence pour
le type de toutes les autres. S'étant dit que les esprits existent, connaissent,
veulent, agissent, et que leurs actions se manifestent par des changements
matériels, il devait chercher en lui quelque chose de semblable : ses
connaissances, ses volontés, le principe de ses actions ont été attribués à une
substance immatérielle, qui a reçu différents noms selon ses opérations. Cette
ébauche contient l'origine de la plupart des idées reprises depuis : la
littérature a gardé les fictions qu'on tenait jadis pour réelles, les sciences
physiques les observations que ces fictions expliquaient, la philosophie y a
puisé ses systèmes et les religions leurs croyances.\footnote{Vue 26, page 16
barrée (verso du f. 12). « les esprits existent, ils connoissent, ils veulent,
ils agissent » est souligné d'un trait ondulé. « devoit » est écrit au-dessus de
mots biffés dont le second est illisible ; « expliquoient » est écrit sur
« expliquer », au-dessus de « avoient servi à », biffé. Un mot biffé après « des
lois immuables » est illisible.}

\subsection*{Un seul être, l'infini, l'éternité}

\pagerange{26}{29}
Les observations se multiplient ; la régularité des mouvements célestes et la
constance des phénomènes terrestres révèlent des lois immuables. Les volontés
d'une foule de personnes n'ont pas ce caractère ; mais un seul homme, chargé de
diriger plusieurs genres d'actions, établirait un ordre constant qui le
dispenserait de l'attention de détail, et la société en offrait des exemples.
L'homme a dit alors : « un seul être a voulu l'univers, et il le gouverne ; ses
volontés sont immuables. » Nous voyons naître nos semblables, nous avons
commencé : l'univers a donc eu un commencement. Nous avons une âme immatérielle,
force motrice de nos actions : l'être des êtres est immatériel, il a tout créé
et agit sur toutes choses. Il a créé l'univers : il existait donc avant lui.
Là, l'esprit arrivait au bout des analogies ; au-delà il n'avait plus de
modèle, donc plus d'idée de l'être ; cette négation d'idée, cette limite de la
pensée a reçu un nom, l'infini, et, pour la durée, l'éternité. Le créateur n'a
pas commencé ; il ne doit pas finir ; il est éternel.\footnote{Vue 27, f. 13r,
page 17 barrée. Trois petits mots au-dessus de « L'état » ne se lisent pas
sûrement. La citation porte des guillemets sur le feuillet.}

Elle s'arrête ici sur sa propre démonstration — « Dans ce qui précède, on ne
voit pas clairement comment on a été conduit à la dernière partie de la
proposition : il ne doit pas finir » — et la complète. Nous voyons des
existences semblables à la nôtre commencer et finir ; en deçà et au-delà, le
temps ne leur appartient pas, mais d'autres existences le remplissent ; ces
limites sont donc relatives, et la durée d'une vie est comprise entre deux
limites de même genre. L'analogie voulait que l'existence dont on avait reculé
l'origine jusqu'à la limite absolue fût elle aussi bornée par deux limites de
même genre : elle ne devait pas finir, puisqu'elle n'avait pas commencé.

L'intelligence s'étant déjà approprié la spiritualité devait prendre aussi
possession de l'éternité. Sa méthode constante est de transporter hors d'elle
les lois de sa propre existence, de chercher dans les analogies ce qui manque à
ses notions, et de reporter ensuite sur elle-même ce que les objets extérieurs
lui ont suggéré : elle était ainsi conduite à l'éternité de l'âme. Cette
opinion a eu de nombreux partisans ; c'est pourtant l'idée, moins analogique,
de la simple immortalité qui a prévalu. De même, quand l'homme eut appris des
arts mécaniques qu'on peut transformer les agents naturels et leur imprimer un
mouvement durable, l'analogie le conduisit à penser que l'être unique qui
gouvernait le monde en était aussi l'architecte.\footnote{Vue 27 (fin) et vue
28, page 18 barrée (verso du f. 13). La syntaxe de la phrase sur la méthode
n'est pas tout à fait arrêtée : « nous l'avons déjà » est porté en marge, suivi
de mots biffés, et deux mots biffés sont illisibles. La phrase « Elle se
trouvoit donc naturellement conduite à l'éternité de l'ame » est une addition
marginale appelée par une croix ; elle remplace « Cette méthode », biffé, et les
mots d'interligne qui l'accompagnent (« du reste », « de l' ») ne sont pas
sûrs.}

Elle note un paradoxe. On avait été conduit directement à dire que le créateur
n'a pas commencé ; qu'il ne doive pas finir en est, pour ainsi dire, le
symétrique. Or, en s'appropriant ce genre de limite, l'homme ne l'adopte plus
pour origine : il en fait le terme de son existence immatérielle, qui aura
commencé et ne finira pas. L'explication, promet-elle, viendra avec l'examen du
lien entre la morale et les croyances.\footnote{Vue 29, f. 14r (le 4 écrit
par-dessus un premier chiffre), page 19 barrée. « Voici » est encadré d'un
trait. Cet examen promis de la morale et des croyances ne se trouve dans aucune
page du volume.}

\subsection*{L'esprit de système, la phrase de d'Alembert, et l'homme au centre}

\pagerange{29}{33}
Voilà, dit-elle, la marche la plus simple : l'intelligence, qui veut pour ses
propres ouvrages la simplicité, l'ordre et les proportions, attribue à la
volonté du créateur celles qu'elle remarque dans l'univers. Mais l'esprit
philosophique ne pouvait se contenter d'une explication qui s'applique
également aux faits les plus contraires, car une volonté est par essence plus
ou moins arbitraire. Il lui fallait un appui plus ferme ; il chercha partout
des lois de nécessité, et y soumit tantôt la toute-puissance divine, tantôt la
matière et ses accidents.\footnote{Vue 29. La construction « par son essence
en plus ou moins arbitraire » est restée telle au manuscrit. Quatre lignes
biffées suivent, reprises plusieurs fois entre les lignes et en partie
illisibles ; on y lit que « l'imagination sans cesse tourmentée » essaya une
infinité de manières de se rendre compte des réalités qui la frappaient.}

Pressé de trouver au dehors des ressemblances avec son modèle intérieur, mais
ignorant encore et la quantité de chaque phénomène et leurs rapports, l'homme
de génie cherchait une dépendance mutuelle entre les faits, les ordonnait
selon ses convenances intellectuelles, et demandait aux analogies ce qui
manquait à ses connaissances positives. L'esprit de système devait égarer :
il mettait à la place des certitudes non acquises mille conjectures hardies,
que des observations nouvelles démentaient ensuite, et qui léguaient aux
générations suivantes de véritables préjugés sur les faits encore inconnus. Il
n'a pourtant jamais cessé d'être guidé par la prévision de la vérité.

On a voulu, en dernier lieu, réunir en un seul faisceau les branches de la
science, et d'Alembert, dans le \emph{Discours préliminaire} de
l'\emph{Encyclopédie} (1751), l'a dit en une phrase qu'elle cite : « l'univers,
pour qui sauroit l'embrasser d'un seul coup d'œil, seroit un fait unique, une
grande vérité ». Ces paroles, écrit-elle, renferment le secret des efforts de
l'esprit humain ; elle y reviendra pour en faire le point de départ de sa
thèse.\footnote{Vue 30, page 20 barrée (verso du f. 14). Le feuillet ne nomme
pas d'Alembert ici, mais « L'auteur de la préface de l'Encyclopédie », qui
« dit, en la terminant » ; il le nomme à la vue 55. Deux précisions sont de
cette lecture. La phrase n'est pas à la fin du \emph{Discours} : elle se trouve
dans sa première partie, là où d'Alembert traite de l'enchaînement des
connaissances et du véritable esprit systématique. Et le texte imprimé ne dit
pas tout à fait ce qu'elle cite : « d'un seul point de vue » et non « d'un seul
coup d'œil », avec « ne serait, s'il est permis de le dire, qu'un fait unique et
une grande vérité ». Elle cite de mémoire, ou abrège. Sur le feuillet,
« l'Encyclopédie », « en un seul faisceau », « sans doute » et quelques autres
additions de cette page sont d'une écriture plus fine et plus droite, peut-être
d'une autre main ; « et » (trouvoit dans les analogies) est incertain ; un mot
biffé devant « l'effet inévitable » est illisible.}

Chaque système voulait ramener les faits connus à un fait unique : on posait
entre eux des relations de cause à effet, on voulait une raison de leur être,
on cherchait l'unité, l'ordre, les proportions, parce que ce sont les
caractères du vrai. Faute de pouvoir les établir, on généralisait, plus ou
moins heureusement, les résultats certains : une vérité connue imprimait son
caractère à un vaste système, et mille suppositions comblaient l'intervalle
entre elle et les vérités secondaires qui semblaient s'y rattacher. Une idée
dominait partout : l'homme s'est cru le modèle de tous les êtres, le but vers
lequel ils tendent, le centre de l'univers. Le soleil, la lune, les étoiles
presque invisibles étaient faits pour fertiliser ses champs et éclairer ses
veilles ; le brin d'herbe, l'animal du désert, le coquillage du fond des mers
avaient leur utilité, c'est-à-dire étaient faits pour lui. Erreur du jugement
d'abord, que l'amour-propre renchérit, et que les religions consacrèrent ; on
en voit encore la trace.\footnote{Vue 31, f. 15r ; en haut à gauche un nombre
barré dont on lit le 2, le second chiffre se confondant avec le trait. Vue 32,
page 22 barrée (verso du f. 15). « on les \emph{concentroit} dans des relations
imaginaires » : la fin du verbe est surchargée, et la lecture incertaine.}

Deux branches de l'ancien savoir, longtemps estimées, sont aujourd'hui tenues
pour fausses : l'alchimie et l'astrologie judiciaire. La première enseignait
que le corps humain est l'abrégé de l'univers ; ses substances portaient les
noms des organes auxquels on les disait semblables — le « foie de soufre » est
encore dans la langue commune. Elle voulait l'unité, elle aussi : elle
cherchait la panacée, remède universel, et l'alcahest, dissolvant universel qui
devait réduire toutes les substances à un seul élément, l'eau ; elle donnait
aux métaux les noms des planètes et établissait entre eux des rapports.
L'astrologie judiciaire, de son côté, jugeait de l'influence des astres sur le
sort de chacun ; l'homme, persuadé de son importance, se croyait menacé par les
comètes, et les grands ne doutaient pas que leur mort ne fût annoncée par ces
astres vagabonds, qui n'auraient pas pris la peine de visiter la terre s'ils
n'avaient été chargés d'avertir ses habitants d'un aussi grand
malheur.\footnote{Vue 32 et haut de la vue 33. Le feuillet écrit « l'alcali ou
le dissolvant universel ». Le dissolvant universel des alchimistes s'appelait
l'\emph{alcahest} (le mot se lit chez Paracelse, et van Helmont en fit la
théorie) ; l'alcali est tout autre chose, une base. La réduction de toute substance à l'eau est bien celle
que van Helmont prêtait à l'alcahest : cette coïncidence n'est pas une
citation, et le feuillet ne nomme personne. Le « foie de soufre »
(\emph{hepar sulphuris}, un mélange de sulfures de potassium) doit son nom à
sa couleur. Le feuillet écrit ensuite « L'astronomie judiciaire », après
« astrologie judiciaire » quelques lignes plus haut ; cette lecture garde le
second terme, qui est le sien. Mots biffés illisibles : quatre, dont trois au-dessus de
« qui » et un devant « ressemblances ».}

\section*{Juger des choses par nos idées : le doute de Kant}

\pagerange{33}{41}
La transition est brusque et marquée sur le feuillet : un « II » romain à
l'encre épaisse sous le folio 16, et, en marge de l'alinéa, « M. » suivi d'un
nom illisible commençant par D, souligné d'un trait qui revient vers le
texte.\footnote{Vue 33, f. 16r, page 23 barrée. Ni le « II » ni le nom ne sont
de sa main, et rien sur le feuillet ne dit ce que le chiffre divise. Ce n'est
pas un titre de chapitre, et cette lecture ne le traite pas comme tel.} Des
erreurs antiques, on passe à celles qu'on commet encore.

\subsection*{L'habitude de juger d'une chose par l'idée qu'on s'en fait}

\pagerange{33}{34}
Nous avons renoncé à l'astrologie, mais nous gardons l'invincible habitude de
juger de la nature des choses par la possibilité de nous en former une idée :
nous affirmons ou nions une proposition selon que nous pouvons ou non concevoir
ce qu'elle énonce. Nous disons que la matière est divisible à l'infini, parce
qu'il est facile de continuer la division arithmétique à l'infini ; et qu'elle
ne peut penser, parce qu'elle est divisible et que l'unité de nos opérations
intellectuelles répugne à la divisibilité. Or nous n'en savons rien, ni
\emph{a posteriori}, puisque l'expérience n'atteint pas ces questions, ni
\emph{a priori}, puisque la matière ne nous est connue que par des perceptions
et que son essence nous échappe.\footnote{Vue 33. Le feuillet écrit « à
posteriori », « à priori », soulignés. Les deux exemples sont exactement ceux
que Kant range parmi les questions que la raison ne peut trancher — la
divisibilité à l'infini dans la deuxième antinomie de la \emph{Critique de la
raison pure} (1781), la simplicité de ce qui pense dans les paralogismes ; la
coïncidence est notée ici, le feuillet ne renvoie pas à Kant à cet endroit.}

On croirait, à notre assurance, que nous avons fait comme le géomètre, qui
définit son objet si exactement que la définition en contient toutes les
propriétés. La différence est entière : au lieu d'une « équation absolue » qui
renferme tout l'objet, nous ne connaissons de la matière que quelques
propriétés relatives à nos sens. Et pourtant, dans des questions qui ont si peu
de prise, nous disons « il est évident », « il est absurde », « il faut être de
mauvaise foi pour ne pas en convenir ». La philosophie a fait des progrès réels,
conclut-elle, mais devra encore beaucoup changer pour atteindre
l'exactitude.\footnote{Vue 34, page 24 barrée (verso du f. 16). Les onze
premières lignes sont d'une encre plus pâle que la suite, les corrections d'une
encre noire. « avons » (à balancer les probabilités) est surchargé et
incertain ; un petit mot barré au-dessus de « lorsque » est illisible.}

Elle rappelle alors sa proposition du premier chapitre : il y a en nous un
sentiment profond d'unité, d'ordre et de proportion qui guide tous nos
jugements ; dans l'ordre moral il donne la règle du bien, dans l'ordre
intellectuel la connaissance du vrai, dans les choses d'agrément le caractère
du beau. Il nous est difficile de savoir si ces conditions résultent
immédiatement des lois de l'être, ou seulement d'un rapport entre toute autre
réalité et notre existence. Des philosophes ont approché la question : les uns
ont cherché dans les causes de nos sensations des qualités qui leur
correspondent, d'autres ont nié l'existence des objets
extérieurs.\footnote{Vue 34. Deux corrections montrent la pensée qui se
retient : « absolument impossible », biffé, devient « difficile » (de savoir),
et « le résultat nécessaire », biffé, devient « le résultat immédiat ». Le
feuillet ne nomme aucun philosophe. Les deux positions répondent, en termes
d'aujourd'hui, à la distinction lockienne des qualités premières et secondes,
et à l'immatérialisme de Berkeley — à ceci près que Berkeley niait l'existence
d'une substance matérielle, non celle des objets tels que nous les percevons ;
« nier l'existence des objets extérieurs » force sa thèse. Le mot du feuillet
est « causes occasionnelles », qui est aussi le nom de la doctrine de
Malebranche ; il paraît pris ici au sens courant.}

\subsection*{La question de Kant, et l'homme qui ne connaîtrait que lui-même}

\pagerange{35}{37}
« De nos jours, Kant a discuté une question de ce genre » : il a fait remarquer
que les arguments les plus concluants peuvent être rapportés soit à des
rapports nécessaires, soit aux formes de notre esprit — pour Kant, les formes de
la sensibilité et les catégories de l'entendement —, de sorte qu'aucune
décision rationnelle parfaite n'est permise, quant au raisonnement \emph{a
priori}. Ce doute est légitime, dit-elle, puisqu'on ne peut comparer aucun
autre jugement à celui de l'homme. Mais l'opinion qui attribue à l'être en
lui-même l'unité, l'ordre et les proportions que nous cherchons partout a pour
elle des inductions, qu'elle va exposer.\footnote{Vue 35, f. 17r, page 25
barrée. Le début est très repris et sa syntaxe n'est pas arrêtée : « sa
remarque expresse » est une lecture incertaine, deux petits mots en marge
(« porte par », « ce que ») sont ramenés par un trait sans place sûre. Une
correction dit l'attitude envers Kant : « a fort bien établi », biffé, devient
« a discuté ». Pour Kant, ce qui revient à « nous » dans la connaissance n'est
pas seulement l'entendement : ce sont les formes de la sensibilité (l'espace et
le temps) et les catégories de l'entendement, conditions de toute expérience
possible ; le feuillet réduit le tout aux « formes de notre entendement ».}

Observation préliminaire : notre logique est faite de règles dictées par la
raison universelle, et elles ne seraient pas moins certaines pour nous si l'on
voulait qu'elles apprennent seulement à former et à reconnaître les jugements
qu'aucun homme de bon sens ne peut contester. Supposons alors la raison
entièrement isolée : un homme à qui aucun objet extérieur à l'esprit ne serait
connu, livré à ses seules pensées et à la société de ses semblables, et qui
réunirait en doctrine les vérités de son être, de ses rapports sociaux, de ses
affections et de ses devoirs. Il trouverait les idées du bien, du vrai et du
beau que nous connaissons ; morale, logique et règles du goût ne changeraient
pas, et la littérature, appauvrie, subsisterait par l'art des récits, la
peinture des passions et l'invention d'une action poétique. Mais la question
de savoir si les rapports entre les parties d'un sujet sont nécessaires en
eux-mêmes ou seulement pour nos formes intellectuelles ne se poserait pas à ses
philosophes ; peut-être n'en comprendraient-ils pas même le sens, puisque,
ne connaissant qu'un seul mode d'existence, ils n'auraient pas la notion
abstraite de l'être. Leur logique serait la nôtre ; leurs opinions dogmatiques,
très différentes.\footnote{Vue 36, page 26 barrée (verso du f. 17). La syntaxe
de « et à celles qu'il qu'il doit aux sociétés humaines » n'est pas arrêtée ;
la phrase modernisée en garde le sens probable. Un « 2 » et un « 1 » en marge
devant « la peinture des passions » et « des sujets » indiquent peut-être un
ordre à rétablir. « pouvoit être » (leur logique pouvait être la nôtre) est
incertain ; un mot biffé après « ne seroient » est illisible.}

\subsection*{La question élargie : le cercle, et l'enfant qui frappe la table}

\pagerange{37}{41}
Elle définit l'objet du doute. La question de Kant sape la réalité absolue de
toutes nos certitudes : elle réduit à des vérités relatives celles mêmes dont
nous avons les démonstrations les plus claires, et atteint surtout ce que nous
croyons des attributs de l'être ; le type qui nous fait distinguer le bien, le
vrai et le beau conviendrait à notre manière de sentir sans avoir hors de nous
de réalité assurée. Kant, ajoute-t-elle, après avoir rejeté la preuve de
l'existence de Dieu tirée de la nécessité d'une cause première, cherche ailleurs
ce qui manque au raisonnement — dans la raison pratique, en fait, dont
l'existence de Dieu devient un postulat ; elle y voit une concession arbitraire, destinée à
protéger son système des formes intellectuelles.\footnote{Vue 37, f. 18r, page
27 barrée ; écriture plus régulière. Ce que dit le feuillet de Kant doit être
précisé. Kant rejette en effet la preuve cosmologique (par la cause première),
mais aussi la preuve ontologique et la preuve physico-théologique. Le feuillet dit qu'il
« cherche dans le sentiment ce qui manque au raisonnement » ; or ce n'est pas au
sentiment qu'il le demande : dans la
\emph{Critique de la raison pratique} (1788), l'existence de Dieu est un
\emph{postulat de la raison pratique}, exigé par la loi morale : une croyance de la raison, non un
sentiment. Correction significative sur le
feuillet : « système de la causalité », biffé, devient « système des formes
intellectuelles ». La construction « à cette de l'univers » est restée telle.}

Sa proposition fondamentale tient : il n'y a qu'un seul modèle du vrai, dont
les copies diffèrent comme les objets qui en reçoivent l'empreinte ; en morale,
en science, en littérature, dans les arts, nous cherchons l'unité d'existence,
l'ordre et les proportions entre les parties d'un même tout. D'où sa question,
plus large que celle de Kant : ce modèle, le devons-nous au fait d'exister,
considéré abstraitement — tout être intelligent trouverait-il en lui les
conditions sans lesquelles aucune existence n'est possible ? — ou au mode
particulier de notre être ? La question de Kant y est comprise : notre logique
est-elle celle de la raison absolue, ou seulement de la raison humaine ?

Kant a remarqué notre tendance à chercher la cause de tout ce qui nous frappe.
Pour elle, cette tendance avertit que nous ne voyons pas l'objet en entier : il
se présente comme une fraction, nous en demandons l'unité ; nous le voyons comme
partie, nous cherchons le tout. Exemple : frappés d'une propriété des sinus et
cosinus, nous demandons pourquoi elle a lieu, car nous n'avons sous les yeux
qu'une partie du sujet ; mais si nous remontons à l'équation du cercle, notre
curiosité est satisfaite : nous avons l'essence, une existence complète, qu'une
intelligence quelconque comprendrait de la même manière.\footnote{Vue 38, page
28 barrée (verso du f. 18), et vue 39, f. 19r, page 29 lue avec quelque doute.
L'exemple doit être précisé. De l'équation $x^2 + y^2 = 1$ du cercle unité, où
l'on pose $x = \cos\theta$ et $y = \sin\theta$, se tire immédiatement la
relation $\cos^2\theta + \sin^2\theta = 1$ ; mais les autres propriétés des
sinus et cosinus — les formules d'addition, par exemple — demandent davantage
que l'équation : la définition de l'angle par la longueur d'arc et la
géométrie des rotations. L'équation ne « renferme » donc pas toutes les
propriétés au sens où le feuillet le dit. Ce que Kant dit de la recherche des
causes répond, dans la \emph{Critique de la raison pure}, à l'exigence de la
raison de remonter à la totalité des conditions, l'inconditionné ; la
réponse du feuillet, « la partie cherche son tout », en est proche, et le
feuillet ne le dit pas.}

Mais de tels sujets sont rares : ils appartiennent aux mathématiques pures,
tandis que la logique s'applique à tout. On a vu que la question de sa
certitude absolue ou relative est insoluble \emph{a priori}, et que l'homme
isolé, s'il la comprenait, n'hésiterait pas à déclarer ses nécessités
intellectuelles absolues. Il irait plus loin, et penserait souvent à l'inverse
de nous. Nous avons établi que la matière ne pense pas ; lui n'aurait pas imaginé
deux substances en lui, rien ne l'aurait forcé à supposer des volontés
invisibles, il n'aurait pas douté de l'unité de son existence ; et si des corps
inertes lui avaient été présentés, il les aurait crus doués de sentiment et de
pensée, jusqu'à ce que l'expérience réforme ce dogme. Ainsi les enfants qu'on a
préservés du contact des objets attribuent au corps qui les heurte la volonté
de les frapper ; la loi du talion, « loi de justice innée », les pousse à
rendre le coup, et la nourrice qui les y invite ne fait que suivre leur désir
d'être vengés. L'observateur dit que l'enfant raisonne mal ; ses idées dérivent
pourtant du même sentiment d'analogie qui a mené l'homme, placé autrement, à des
dogmes opposés.\footnote{Vue 39, en marge, d'une grande écriture, « Pennequin »,
souligné ; plus bas, « Absolue » en marge, mis par la transcription après
« certitude ». « il hésiteroit pas », sans « n' », est ainsi au manuscrit. Le
verbe avant « pas douté » est surchargé (« n'eût », incertain), et, dans « il
n' … pu manquer », un mot biffé et un mot surchargé ne se lisent pas ;
deux autres mots biffés sont illisibles. Vue 40, page 30 barrée (verso du
f. 19) : « talion » est écrit d'une plume chargée et lu avec doute.}

Où trouver la solution ? Les raisonnements \emph{a priori} ne l'atteignent pas,
puisqu'ils sont l'œuvre de la raison même qu'on veut juger ; les preuves
extérieures non plus, puisque l'analogie conduit chaque observateur, selon sa
position, aux opinions les plus contraires. Si les conjectures de l'homme
avaient toujours été vérifiées, si l'observation n'avait jamais démenti ses
décisions théoriques, qui douterait de l'absolu de nos nécessités logiques ?
Les formes de l'intelligence auraient-elles le pouvoir de plier les objets à
leurs convenances ?\footnote{Fin de la vue 40 et vue 41, f. 20r, page 31
barrée. Le feuillet écrit « l'absolutisme de nos nécessités logiques » ; le mot
« absolutisme », qu'elle emploie encore aux vues 43, 53 et 54, a ici le sens de
« caractère absolu » et n'a rien de politique. « avait » est écrit ainsi à la
vue 40, parmi des « avoit ».}

\section*{Deux questions, et la réponse de l'histoire}

\pagerange{41}{55}
Nous sommes loin de cette position heureuse : l'histoire des sciences montre
mille écarts, et l'esprit a dépensé plus d'efforts à détruire ses ouvrages qu'à
en construire. Chaque système rendait compte des faits à l'aide de suppositions
hasardées, devenait insuffisant, et son autorité faisait alors obstacle à la
vérité. D'où deux questions, qu'elle pose et traite dans l'ordre. Si nous
portons en nous le type du vrai, pourquoi tant de méprises ? Si notre logique
est le recueil des préceptes de la raison absolue, comment avons-nous pu errer
si longtemps parmi les suppositions gratuites ?\footnote{Vue 41. « Le systèmes
satisfesoit » est ainsi au manuscrit (« Le » écrit sur « Chaque », biffé).
De la vue 41 à la vue 60, la pagination à l'encre court de 31 à 50 ; seuls
31, 37, 41, 43, 46 et 50 (vues 41, 47, 51, 53, 56, 60) se lisent avec
assurance, les autres sont serrés sous le trait et lus dans la suite. 
Au-dessus de « Si », biffé, une addition de deux signes soulignée, dont le premier
n'est pas lu.}

\subsection*{Pourquoi tant d'erreurs, et comment une logique sûre peut égarer}

\pagerange{41}{44}
À la première : le type du vrai ne s'est jamais tu au milieu des écarts. Chaque
système a été inspiré par une vérité incontestable ; l'homme de génie, frappé
de son importance et persuadé de l'unité de l'être, a voulu tout y rapporter,
et a rempli d'hypothèses les intervalles entre les points solidement établis.
Mais il n'a jamais confondu, dans sa conscience, la certitude qui fondait son
œuvre avec la probabilité, souvent faible à ses propres yeux, des suppositions
qui en liaient les parties. C'est donc dans la pensée des inventeurs qu'il faut
étudier l'intelligence, et l'histoire du genre humain se réduit, à cet égard, à
celle d'un petit nombre d'hommes. Les esprits communs, juges éclairés dans le
cercle de leurs intérêts, ne voient pas au-delà. Parées des formes d'une
imagination élevée, les doctrines ont été adoptées avec enthousiasme et
enseignées dans les écoles, où il s'agissait de les savoir, non de les juger ;
on les commentait de mille manières que le maître n'eût pas admises ; la
simplicité première se perdait, les erreurs s'accumulaient, et l'enseignement y
restait attaché jusqu'à ce qu'une hypothèse plus conforme aux observations
vînt satisfaire ce besoin de savoir qui a si longtemps devancé la science
véritable.\footnote{Vue 42, et vue 43, f. 21r, page 33 ; à côté du folio 21, un
« 8 » biffé à l'encre. « on » (faisoit sur ses écrits mille commentaires) est
souligné, surchargé et incertain.}

À la seconde : un faux raisonnement que la raison humaine peut reconnaître tel
ne prouve rien contre le caractère absolu de nos nécessités logiques. Reste le
cas de l'homme de bonne foi et de jugement sûr qui, parti d'un principe certain
et raisonnant avec exactitude, arrive à des conclusions que les faits démentent.
La cause est dans la définition : il est extrêmement difficile d'énoncer le
principe assez précisément pour être sûr qu'il y est contenu tout entier, et
qu'aucune idée étrangère ne s'y est glissée. Et cette difficulté tient aux
langues. Nées des échanges ordinaires, elles sont exactes pour les choses
présentes ou bien connues ; pour les sujets philosophiques il a fallu prendre au
figuré des mots clairs au propre, dont la précision ne survivait pas au
changement de sens. On a cru y parer en forgeant des mots nouveaux ; mais ils
demandaient à être définis à leur tour, et les termes techniques, interprétés
diversement par ceux qui cherchaient dans un système accrédité l'appui de leurs
idées, sont devenus l'une des sources les plus fécondes de la divagation
philosophique : mêmes mots, opinions disparates.\footnote{Vue 44, page 34
(verso du f. 21). « l'appui » : l'article est incertain. Deux mots biffés sont
illisibles, après « dont nous » et devant « Les ».}

\subsection*{Des nombres à Newton}

\pagerange{45}{48}
La digression historique commence par les mathématiques. Dans un temps très
reculé, les propriétés générales des nombres, des figures simples et des corps
réguliers ont attiré les esprits nés pour les sciences exactes. Les idées y sont
d'une extrême simplicité ; on avait les figures sous les yeux et l'on ne pouvait
leur attribuer des propriétés qu'elles n'avaient pas. Des remarques
multipliées ont donné une foule de théorèmes, énoncés avec une précision
entière — dans l'Antiquité en langue ordinaire, sur des figures désignées par
des lettres, car les signes de l'algèbre ne viendront qu'avec Viète et
Descartes. Dès leur naissance, les mathématiques ont offert à l'esprit la
réalisation entière du type du vrai, qu'il cherchait en vain ailleurs, où mille
suppositions s'étaient incorporées à quelques vérités ; et l'homme d'esprit
juste sentait, sous les formes absolues de l'enseignement philosophique, que
l'étude ne le mènerait à aucune certitude.\footnote{Vue 45, f. 22r, page 35. Le
feuillet dit que « quelques signes dont la signification ne pouvait être
équivoque ont suffi pour représenter avec précision des idées d'une exactitude
parfaite ». Pour les Grecs, c'est inexact : leur géométrie et leur arithmétique
s'écrivent en phrases, les figures désignées par des lettres, et le
symbolisme algébrique est une création des XVI\textsuperscript{e} et
XVII\textsuperscript{e} siècles (Viète, \emph{In artem analyticen isagoge},
1591 ; Descartes, \emph{La Géométrie}, 1637). Un mot en paraphe sous la fin du
paragraphe est illisible.}

Aussi, quand les sciences physiques, morales, religieuses et politiques étaient
encore chargées de doctrines hypothétiques, la géométrie inspirait-elle un
enthousiasme qu'on ne retrouve plus : comment l'homme sensible aux conditions du
vrai, sortant des systèmes obscurcis par des commentateurs qui se
contredisaient de cent manières, n'aurait-il pas été transporté de trouver la
vérité au plus haut degré d'évidence ?\footnote{Vue 46, page 36 (verso du
f. 22). « évidence », biffé, devient « clairvoyance » sur le feuillet.}

Descartes osa douter publiquement des doctrines de l'école. Il ne renonça pas
pour autant à l'espoir, tant de fois déçu, de réaliser la copie fidèle du type
de l'être, et reconstruisit l'univers sur un plan nouveau ; mais l'exemple qu'il
avait donné servit bientôt à faire rejeter son propre système. Il rendit ainsi à
la raison un service immense et lui donna son indépendance ; l'hypothèse des
tourbillons, qui appartenait encore à l'âge finissant, en marqua la dernière
limite, et les efforts de l'esprit changèrent de direction. Surtout, il réunit
les deux parties jusque-là distinctes des mathématiques par l'application de
l'algèbre à la géométrie — Fermat y parvenait au même moment de son côté —, et ouvrit une route où l'on ne pouvait plus
s'égarer ; la langue du calcul, pliée à cet usage, devint bientôt capable
d'exprimer les grands faits du ciel.\footnote{Vue 46 (« Descartes » est écrit
sur un « Il » ; biffé : « ce type étoit empreint avec force dans son génie
créateur ») et vue 47, f. 23r, page 37. Le feuillet attribue la réunion de
l'algèbre et de la géométrie à Descartes seul. Fermat l'a faite de son côté, au
même moment — son \emph{Ad locos planos et solidos isagoge} circulait vers 1636
et ne fut imprimé qu'en 1679 —, et l'« art analytique » de Viète appliquait déjà
l'algèbre aux problèmes de géométrie. Un mot court au-dessus de « fut », biffé,
est illisible.}

Newton vint, armé d'un nouveau calcul, et l'unité, l'ordre et les proportions de
l'univers, que le sentiment du vrai faisait chercher depuis si longtemps,
devinrent des vérités mathématiques. Il établit la loi des mouvements célestes
et, par une analyse très fine, les mesura ; il fit de l'optique une science
nouvelle, et devina sur les corps réfringents des vérités que la chimie ne
vérifia que longtemps après : ayant remarqué que le diamant réfracte la lumière
comme les corps huileux et inflammables, il le conjectura combustible, ce que
les expériences de combustion du diamant (Lavoisier et d'autres, 1772) et
Tennant (1797, le diamant est du carbone) ont confirmé.\footnote{Vue 47. Trois
points du feuillet appellent une précision. « Armé d'un nouveau genre de
calcul » : Newton a bien inventé le calcul des fluxions (vers 1665--1666),
Leibniz de son côté le calcul différentiel (publié en 1684) ; mais les
\emph{Principia} (1687) exposent leurs démonstrations presque entièrement sous
forme géométrique, par les « premières et dernières raisons » ; la mise en
équations différentielles de la mécanique céleste est l'œuvre de ses
successeurs, Euler, Clairaut, d'Alembert, Lagrange, Laplace. « Son génie avoit
reconnu la cause des mouvemens célestes » : Newton a établi la loi de la
gravitation, et s'est expressément refusé à en assigner la cause (\emph{hypotheses
non fingo}, scholie générale de 1713) ; cette lecture écrit « la loi ». Enfin le
feuillet ne précise pas quelles vérités sur les corps réfringents ; l'exemple du
diamant est de cette lecture, comme le plus connu de ceux que l'\emph{Optique}
(1704) propose.}

C'est de Newton qu'elle date l'alliance des mathématiques et de la physique.
Les Anciens, dit-elle, connaissaient déjà la pratique et la théorie de la
mécanique et de l'hydrodynamique, comme en témoignent de grands ouvrages et les
livres d'Archimède, et l'idée de quantité est si liée à celle de force que ces
deux sciences sont essentiellement mathématiques. Ce que les Anciens
possédaient en fait, c'est la statique et l'hydrostatique : l'équilibre du
levier et des corps plans, les corps flottants ; la dynamique — le mouvement
sous l'action des forces — est de Galilée et de Newton, et l'hydrodynamique de
Daniel Bernoulli (1738) et d'Euler (1757). Leurs éléments, avec ceux de
l'algèbre et de la géométrie, formaient tout le domaine des idées exactes ;
ailleurs, les vains efforts du génie et les erreurs sans nombre que traînaient
les doctrines des premiers inventeurs.\footnote{Vues 47 et 48 (page 38, verso
du f. 23). Le feuillet écrit : « La mécanique et l'hydrodynamique avoient été
connues des anciens. De grands travaux et les livres d'Archimède attestent
qu'ils en savoient et la pratique et la théorie. » Les livres d'Archimède sont
\emph{De l'équilibre des plans} et \emph{Des corps flottants} : de la statique
et de l'hydrostatique, non de la dynamique ni de l'hydrodynamique.}

Le langage mystérieux des philosophes, plus obscur que leurs idées, contrastait
singulièrement avec la langue précise des sciences exactes. Tant que les
géomètres furent peu nombreux et isolés, eux seuls le voyaient, et en tiraient
un profond mépris pour les autres sciences. Quand les phénomènes célestes se
rangèrent sous les lois du calcul, l'étude des mathématiques se répandit, et les
bons esprits furent frappés d'une manière d'argumenter si différente de celle
de l'école. L'astronomie physique remplaçait des hypothèses discréditées ; cette
révolution subite ébranla les préjugés et alarma ceux qui avaient intérêt à leur
règne : ils redoutaient jusqu'aux vérités les plus étrangères à leurs doctrines,
et, comme le peintre qui écarte du regard tout objet réel, une sorte de
perspective morale les avertissait du danger des comparaisons.\footnote{Vue 48
et haut de la vue 49 (f. 24r, page 39). « L'astronomie » (physique) est une
lecture incertaine. Trois lignes biffées, qui portaient en interligne
« supériorité », disaient que « le contraste entre la rigueur de leur méthode et
l'argumentation stérile des école » frappa tous les esprits ; la reprise ne
parle plus que d'« une manière d'argumenter si différente ». Un mot en paraphe,
souligné, suit la phrase sur la perspective morale ; il est illisible.}

\subsection*{Bacon, le « comment » et le « combien »}

\pagerange{49}{51}
Tandis que le système du monde offrait le spectacle d'un mécanisme simple dans
son principe et fécond dans ses conséquences, la physique terrestre restait
chargée de suppositions. L'attachement aux routines ne pouvait tenir longtemps
contre l'espoir de succès semblables dans d'autres études : on sentit qu'il
fallait tout refaire. Elle place ici le conseil de Bacon — qui l'avait en fait
donné bien avant, son \emph{Novum Organum} (1620) précédant le
\emph{Discours de la méthode} (1637) et les \emph{Principia} (1687).\footnote{Vue
49. Le feuillet écrit, après Newton : « on sentit qu'il falloit tout refaire.
Bacon en donna le conseil. » L'ordre du récit fait suivre le conseil de
l'exemple ; l'ordre des dates est inverse. Mots illisibles : une première
correction biffée au-dessus de « l'envie d'imposer ».}

Les Anciens, guidés par la métaphysique, avaient peu observé, comme s'ils
craignaient que les faits ne démentissent leurs systèmes ; à la Renaissance, on
admira à juste titre leur littérature, puis tous leurs ouvrages également, on
adopta leurs idées, et les controverses ne portèrent plus que sur la manière de
les interpréter. Jusque-là on cherchait les causes ; on se mit alors à
considérer les phénomènes en eux-mêmes : au lieu du \emph{pourquoi}, on voulut
savoir le \emph{comment}. Des observateurs laborieux renoncèrent pour eux-mêmes
à expliquer, afin de léguer à leurs successeurs une masse de connaissances
positives dont la liaison se montrerait plus tard. « C'est de cette époque que
datte la connoissance de la nature. Jusque là l'homme l'avoit imaginée ; il la
vit alors pour la première fois. » On chercha ensuite à mesurer tout ce qui est
mesurable, et au \emph{comment} se joignit le \emph{combien} ; les phénomènes
devinrent calculables, et les plus simples furent l'objet des recherches des
successeurs de Newton.\footnote{Vue 49 (« on \emph{diroit} » et « le démenti
\emph{à} », lectures incertaines) et vue 50, page 40 (verso du f. 24) ; les
mots \emph{pourquoi}, \emph{comment}, \emph{combien} sont soulignés sur le
feuillet. Un signe barré et taché d'encre, en marge, est illisible. « Les
anciens \ldots{} avoient peu observé » est un jugement trop large :
l'astronomie d'Hipparque et de Ptolémée, les ouvrages biologiques d'Aristote,
la médecine de Galien reposent sur des observations nombreuses.}

De nos jours, poursuit-elle, l'esprit mathématique a tant progressé que la
« physique particulière » — l'étude des phénomènes naturels qui ne relèvent pas
de l'histoire naturelle — a presque disparu en devenant une branche des
sciences exactes, et que la langue du calcul, enrichie de méthodes nouvelles,
traite des questions qu'on lui croyait étrangères. C'était vrai d'une partie de
la physique de son temps : la mécanique, l'optique, la chaleur, l'électricité
statique et l'élasticité venaient d'être mises en équations ; ce ne l'était ni de
la chimie, ni de la plus grande part de l'électricité et du
magnétisme.\footnote{Vue 50. Le feuillet dit que la physique particulière « a
pour ainsi dire disparu, et s'est trouvée transformée en une des branches les
plus importantes des sciences exactes ». La formule est trop forte pour quelque
date que ce soit du premier tiers du XIX\textsuperscript{e} siècle.} Il y a
moins d'un siècle, les personnes les plus instruites tenaient l'algèbre pour un
langage barbare et indéchiffrable ; aujourd'hui ses éléments entrent dans
l'éducation, et la raison publique y a puisé des forces nouvelles.\footnote{Vue
51, f. 25r, page 41. Un mot biffé après « confus » est illisible.}

\subsection*{Les réponses, et la langue du calcul}

\pagerange{51}{55}
Reprenant ses deux questions, elle répond d'abord à la première. Le type du vrai
est fait d'idées abstraites : il nous avertit de ce qui répugne, de ce qui ne
peut exister ensemble, mais il ne suffit pas à nous faire connaître des réalités
particulières. Nous savons que chaque chose a une essence, qui est l'unité du
sujet, qu'elle se divise, qu'il y a de l'ordre et des proportions entre ses
parties ; mais quelle essence, quel ordre, quelles proportions dans un cas donné,
le modèle de l'être ne le dit pas. L'homme de génie qui a suppléé par des
suppositions à ce qui lui manquait d'observations a donc établi des relations
fantastiques, et, si l'on demandait combien de telles conjectures pouvaient se
trouver vérifiées, on trouverait la très petite fraction des succès obtenus
avant l'époque où les observations précises ont remplacé les
assertions.\footnote{Vue 52, page 42 (verso du f. 25), et vue 53, f. 26r, page
43. « le modèle de l'être ne suffit pour nous en informer » : le « pas » manque
au manuscrit, et le sens l'impose. Le feuillet dit qu'on demanderait au
géomètre « combien de fois la théorie des probabilités veut que des conjectures
ainsi formées se réalisent » ; aucune théorie des probabilités ne donne ce
nombre, faute d'un modèle des conjectures et de leurs chances, et la phrase est
une figure. Le mot « veut », en fin de page, paraît biffé ; dans une
première rédaction biffée, « l'analyse de », un mot illisible, « la théorie
des probabilités ».}

À la seconde question, elle a répondu par l'imperfection des langues, qui
introduit à l'improviste dans la définition des idées étrangères au sujet. Cette
explication est aujourd'hui pleinement justifiée par les théories
mathématiques, et le caractère absolu des nécessités logiques ne lui semble plus
pouvoir être mis en doute. Car un petit nombre d'hommes a formé une langue
précise, où la moindre erreur deviendrait sensible : c'est la langue de la
raison dans toute sa pureté ; elle interdit la divagation, signale l'erreur
involontaire, et il faudrait ne pas la connaître pour vouloir la faire servir à
l'imposture.\footnote{Vue 53. Le paragraphe précédent, sur le même sujet, est
annulé d'un long trait oblique et récrit dans celui-ci. La phrase suivante,
commencée au bas de la vue 53 (« et si quelque ») et achevée en tête de la vue
54 (page 44, verso du f. 26), est annulée d'une croix : « si quelquefois il est
arrivé qu'un géomètre ait voulu, dans un intérêt d'amour propre, s'approprier
des inventions qui ne lui appartenoient pas, il n'a pû même en copiant le
véritable auteur dissimuler longtems, que l'idée principale ne lui appartenoit
pas ». La remarque sur le plagiat en géométrie a été retirée ; elle n'est pas
dans la lecture. « même », au-dessus de « quelque », est incertain.}

Cette langue reproduit toutes les conséquences du principe qu'on lui a confié.
Elle peut donc servir à prouver que l'unité d'essence, l'ordre et les
proportions que l'esprit cherche partout n'expriment pas seulement les
conditions de notre satisfaction intellectuelle, mais appartiennent à l'être,
à la vérité. Lorsqu'on a su rendre une question mathématique — en saisir
l'essence assez simplement pour que l'analyse s'en empare —, la nature vient
sanctionner les oracles de la science : un fait connu s'est présenté comme la
conséquence d'un ordre de choses inconnu, le savant a défini cet ordre, et
l'expérience, abondante en circonstances nouvelles, montre bientôt et le génie
qui l'a deviné et l'excellence de la méthode. Comment douter que le type de
l'être ait une réalité absolue, quand on voit la langue du calcul faire jaillir
d'une seule réalité toutes celles qui lui sont liées par une essence commune, et
l'observation, par une voie toute différente, montrer hors de la pensée un
édifice semblable à celui dont l'homme trouve le modèle en
lui-même ?\footnote{Vues 54 et 55 (f. 27r ; à côté du 27, un second « 27 »
biffé à l'encre). Une correction de la vue 54 montre la pensée qui change :
elle avait d'abord écrit que ces conditions « ne sont pas seulement un besoin
dû à ses formes intellectuelles » (« pas seulement » servant aux deux
rédactions), puis a remplacé « un besoin dû à ses formes
intellectuelles » par « les conditions de notre satisfaction intellectuelle ».
Le vocabulaire kantien des « formes » est retiré à cet endroit. À la vue 55, un
petit « 2 » au-dessus de « lui », et à la suite de la ligne, en plus grands
caractères et souligné, un mot lu avec doute « Deracke ».} L'argument est
celui qu'on nommerait aujourd'hui l'efficacité des mathématiques dans les
sciences de la nature ; il montre que le monde se laisse décrire par nos
théories, non, à lui seul, que nos formes de pensée sont celles de
l'être.\footnote{Remarque de cette lecture. L'expression « déraisonnable
efficacité des mathématiques » est d'Eugene Wigner (1960) ; la question est la
même, et le feuillet ne la tranche que par l'induction qu'elle annonçait à la
vue 35.}

\section*{La thèse : le fait unique est nécessaire}

\pagerange{55}{59}
« Les \emph{préliminaires} qu'on vient de lire nous ont paru nécessaires pour
bien entendre les idées que nous allons exposer » : ils en fixeront le sens et
feront peut-être pardonner ce qu'elles ont de hardi et de nouveau. Nous
avançons, dit-elle, vers la réalisation de ce qui ne fut qu'un pressentiment
chez les auteurs de tant de systèmes prématurés : ramener toutes choses à une
seule, trouver l'unité de l'être. D'Alembert l'a exprimé par la phrase « déjà
citée » : « L'univers, pour qui sauroit l'embrasser d'un seul coup d'œil seroit
un fait unique, une grande vérité. » Elle ajoute : ce fait unique doit être
nécessaire.\footnote{Vue 55. Le mot \emph{préliminaires} est souligné sur le
feuillet. « d'un seul coup d'œil » est ajouté au-dessus de la ligne, et le
dernier mot est couvert d'une tache ; chaque ligne de la citation s'ouvre par un
guillemet en marge. Sur l'écart entre cette citation et le texte de
d'Alembert, voir la note sur la vue 30.}

\subsection*{Essence et nécessité}

\pagerange{55}{56}
Chercher l'essence d'une chose et chercher sa nécessité, c'est tout un : quand
nous connaissons l'essence, nous voyons que l'être qui la possède ne saurait ni
ne pas être, ni être autrement qu'il n'est, et l'esprit, appuyé sur la
nécessité, trouve le repos. L'attrait des sciences exactes n'a pas d'autre
cause : leurs objets sont connus dans leur essence, et l'on ne saurait concevoir
qu'ils n'existent pas ; l'esprit y entre dans la possession de l'être nécessaire,
de la vérité pure. Partout ailleurs nous ne voyons que des êtres dépendants et
des vérités partielles, dont nous cherchons l'origine, et nous admettons sans
peine qu'ils pourraient ne pas exister, ou être autres. Cette disposition tient
uniquement à notre ignorance. Aussi les progrès des sciences, en liant des faits
qu'on croyait isolés, nous forcent-ils, dès que les uns sont établis, à tenir les
autres pour nécessaires : nous les voyons alors comme les parties d'une même
existence, et non plus comme des unités séparées.\footnote{Vue 56, page 46
(verso du f. 27), et haut de la vue 57. « l'origine » est écrit sur « la
cause », biffée ; le mot biffé après « dont » est lu avec doute.}

\subsection*{Contingent, probable, nécessaire}

\pagerange{57}{58}
Pour les faits éventuels, l'analyse calcule la probabilité que l'un arrive
plutôt que l'autre. La réponse est empirique : sans s'occuper des circonstances
qui produisent l'événement, le géomètre dit qu'une cause le fait arriver
quelquefois, et que la probabilité qu'elle l'amène est telle fraction. Réponse
utile, mais qui atteste notre ignorance. Qu'on demande la probabilité qu'une
machine s'arrête à un instant donné : si l'on connaissait parfaitement la force
motrice, les frottements et les résistances, on saurait que l'arrêt à cet
instant est certain ou impossible — l'instant de l'arrêt serait déterminé
exactement. Pour des événements plus compliqués, nous ne savons pas même quelles
connaissances nous manquent ; mais faut-il croire les circonstances
déterminantes arbitraires, sans ordre, étrangères aux conditions de toutes les
réalités connues ? Elle conclut que la distinction du contingent et du
nécessaire est, au fond, celle des faits dont on ignore et de ceux dont on
connaît la nature.\footnote{Vue 57, f. 28r, page 47, et vue 58, page 48 (verso
du f. 28) ; un crochet ouvre la conclusion. Deux mots biffés après « telle
fraction » et un mot effacé d'une traînée d'encre en tête de la vue 58 sont
illisibles. Le feuillet dit que l'impossibilité de l'arrêt « seroit évidente
pour l'instant qui précède immédiatement celui pour lequel la réalisation du
même événement seroit aussi évidente » ; en temps continu il n'y a pas
d'instant immédiatement précédent, et ce que la phrase veut dire est que
l'instant de l'arrêt serait connu exactement. La thèse — la probabilité mesure
notre ignorance d'un monde déterminé — est celle que Laplace a formulée dans
l'\emph{Essai philosophique sur les probabilités} (1814) ; le feuillet ne le
nomme pas. La physique postérieure l'a nuancée de deux côtés : la mécanique
quantique pose des probabilités qui ne se ramènent pas à une ignorance
(interprétation de Born, 1926), et le chaos déterministe (Poincaré, 1890 ;
Lorenz, 1963) montre qu'un système parfaitement déterminé peut rester
imprévisible en pratique.}

L'univers, ce fait unique qui tourmente depuis si longtemps les philosophes,
paraîtrait nécessaire s'il était mieux connu ; « on sait que cette opinion a été
soutenue ».\footnote{Vue 58. L'opinion n'est pas attribuée sur le feuillet. Elle
coïncide avec le nécessitarisme de Spinoza (\emph{Éthique}, I, prop. 29 et 33 :
rien n'est contingent, et les choses n'ont pu être produites autrement).}

\subsection*{Dieu, la liberté, la balance}

\pagerange{58}{59}
Les distinctions entre l'intelligence et la matière, dont elle a montré
l'origine, ont repoussé la nécessité jusqu'à Dieu, et l'idée de Dieu a été
formée sur le modèle de notre intelligence. On a dit : Dieu est nécessaire, sa
volonté est libre, il a voulu l'univers. Mais dire que sa volonté est libre,
c'est rompre la chaîne : s'il a pu ne pas vouloir l'univers, l'univers n'émane
plus de lui comme les vérités secondaires émanent de l'unité nécessaire dont
elles font partie. Le modèle suivi est le sentiment de liberté qui accompagne nos
décisions ; or ce sentiment ne nous empêche pas de reconnaître que notre volonté
est souvent entraînée par la nécessité. Nous délibérons réellement, mais nous
nous décidons ; comme une balance chargée des deux côtés, nous oscillons, et le
poids le plus fort fixe la position où le système demeure en repos. La
délibération nous donne le sentiment de la liberté et nous distrait de prévoir
une décision qui, bien que nécessaire, semble avoir été sur le point de
basculer ; aussi une personne qui connaît la situation et le caractère d'une
autre prévoit-elle avec certitude le parti qu'elle prendra, tandis que
celle-ci, étonnée, assure en toute sincérité qu'il s'en est peu fallu qu'elle
n'agît autrement.\footnote{Vue 58 (« nous délibérions » est ainsi au manuscrit ;
« demeure » est écrit au-dessus de « est », biffé) et vue 59, f. 29r, page 49
(« Aussi voyons que une personne » est ainsi au manuscrit). La position qu'elle
défend — une volonté déterminée par le motif le plus fort, et un sentiment de
liberté qui naît de la délibération — est ce qu'on appelle aujourd'hui un
déterminisme psychologique ; le feuillet ne se réclame d'aucun auteur.}

Plus on réfléchit, plus on reconnaît que la nécessité gouverne le monde : à
chaque progrès des sciences, ce qui passait pour contingent se révèle nécessaire,
des liaisons apparaissent entre des branches crues séparées, des lois là où l'on
ne voyait que des accidents. Nous approchons de l'unité de l'être, rêve de
l'Antiquité, dont le modèle est dans le sentiment de notre propre existence.
Reste à fixer son opinion sur ce modèle du vrai, qui a souvent égaré la raison
et qui, dans les temps modernes, la guide si heureusement que ses progrès,
d'abord réservés à quelques savants, se répandent dans toutes les classes,
éclairent les sciences morales et politiques, la physique, les arts chimiques
et mécaniques, et peuvent encore donner aux lettres et aux beaux-arts des
lumières nouvelles.\footnote{Vue 59. « Tâchons » (enfin de fixer notre opinion)
est récrit sur un autre mot et incertain. « ses progrès \ldots{} se répand » est
ainsi au manuscrit.}

\section*{L'être, l'analogie et l'esprit de système}

\pagerange{60}{63}
Une page nouvelle, et, en tête, au crayon et d'une autre main, « Ch 4 » et le mot
« type » ; ce qui suit reprend en effet la question du type.\footnote{Vue 60,
page 50 barrée (verso du f. 29). La ligne au crayon est très pâle ; seuls
« Ch 4 » et « type » y sont lus, le reste est illisible. La marque n'est pas de
sa main, et l'on ne sait pas où, pour la main qui l'a écrite, commençait un
chapitre 3.}

L'homme qui n'aurait d'autre sujet d'étendue que lui-même connaîtrait l'étendue,
et ne douterait pas sérieusement de cette propriété de la matière ; mais ce
qu'il connaîtrait avec la dernière évidence, c'est sa propre existence. L'esprit
humain a essayé du scepticisme, il a pu soutenir que tout ce qui est hors de
l'homme n'est qu'apparence ; mais « le fameux argument, “je pense, donc je
suis” » a ramené la raison à la réalité de son être. Le sentiment de l'être est
celui de la vérité ; il est inséparable de notre existence et précède toute autre
idée. Le bon et le beau dérivent du vrai, mais leur connaissance exige des
comparaisons. Selon qu'on a été frappé de l'une ou de l'autre moitié de cette
proposition, on a soutenu que nos idées sont innées ou qu'elles viennent des
sensations ; les deux opinions sont vraies dans ces limites. Le type du vrai,
nous l'apportons en naissant ; en ce sens l'homme est l'abrégé de l'univers,
puisque l'être, partout où il se trouve, remplit certaines conditions que
l'attention découvrira dans tous les objets réels. Mais cette ressemblance
abstraite est très éloignée de celle qu'on a cherchée jadis, et elle explique
l'erreur qui a séduit l'esprit humain — celle de l'alchimie, qui voyait dans le
corps humain l'abrégé du monde.\footnote{Vue 60 et haut de la vue 61 (f. 30r).
« L'homme, peut-il pas d'autre sujet d'étendue que lui même, connoîtroit
l'étendue » est laissé tel au manuscrit (« peut-il » au-dessus de « s'il
n'avoit », biffé) ; la phrase modernisée suit la première rédaction, qui se
construit. Deux mots biffés, sous « vraies » et devant « dans tous les objets »,
sont illisibles. La phrase finale s'arrête au bas de la vue 60 (« la cause d'une
erreur ») et s'achève en tête de la vue 61 (« qui a séduit autrefois l'esprit
humain ») ; que cette erreur soit celle de la vue 32 est une lecture, appuyée
sur la reprise de la formule « l'abrégé de l'univers ». « Je pense, donc je
suis » est de Descartes, \emph{Discours de la méthode} (1637), IV\textsuperscript{e}
partie ; les idées innées sont la thèse de Descartes et de Leibniz, l'origine
des idées dans les sensations celle de Locke et de Condillac ; le feuillet ne
nomme personne.}

De toutes nos idées, la plus abstraite est celle de l'être, car celle du néant
est toute négative. L'être nous appartient, il éclaire notre intelligence. Le
beau et le bon sont plus composés : nous les devons à la comparaison entre nos
connaissances acquises et notre modèle intérieur. Le grand et le petit, le fort
et le faible, la beauté ou la bonté relatives sont des comparaisons dues aux
sensations et à la réflexion, qu'il serait absurde de croire innées. C'est à
l'uniformité des conditions de l'être qu'il faut attribuer le sentiment
d'analogie qui dirige toutes les opérations de l'entendement, et ce sentiment a
produit des erreurs grossières aussi bien que des pensées heureuses. Comment une
cause constante a-t-elle des effets si différents ? Par les manières diverses
dont on a voulu réaliser ce qu'elle indique. Les conditions de l'être sont
abstraites — sinon elles ne seraient pas universelles. L'esprit de système
consistait à prendre une vérité particulière pour base de tout un ordre de
faits, à ne plus considérer ceux-ci que par leurs rapports, vrais ou supposés,
avec le premier, et à les dépouiller de leur réalité propre pour les revêtir de
celle de la vérité dominante. Au lieu d'analogies, on voulait des identités,
plus simples et donc plus satisfaisantes ; on croyait réaliser le type du vrai
au gré de l'imagination.\footnote{Vue 61, f. 30r. En tête, au-dessus de la
première ligne, un mot souligné lu avec doute « Varlé », un signe en forme de
« § » et une accolade. Un mot biffé devant « le sentiment d'analogie » est
illisible. Vue 62 (verso du f. 30 ; le numéro de page, biffé, ne se lit
pas) : « Le type
\emph{du} vrai », article incertain.}

On devait s'égarer, et pourtant les erreurs ne se sont pas multipliées autant que
le vice de la méthode le ferait croire : elles se sont toutes rapprochées de
certaines vérités, parce que le sentiment du vrai n'a jamais abandonné les
auteurs des systèmes et a retenu leur imagination dans des limites. À l'esprit de
système succède aujourd'hui la recherche méthodique. On comprend mieux la
généralité des conditions de l'être, on cherche comment elles se réalisent dans
chaque sujet, sans plus les confondre avec les conditions particulières de la
vérité dont la découverte fortuite a révélé un ordre de phénomènes. On consulte
l'expérience ; on multiplie les faits en variant les circonstances ; averti par
le sentiment de l'analogie de l'existence de lois qu'il n'aperçoit pas encore,
l'observateur sépare les circonstances qui compliquent les résultats et cherche
les plus grandes et les plus petites influences. Les faits se classent, un ordre
paraît, les lois prévues se manifestent, et une branche nouvelle s'ajoute à la
science. Ce n'en est encore que la partie expérimentale. La théorie naît quand
les faits se prêtent à une expression analytique dont on tire des conséquences
conformes à l'expérience ; les formules révèlent ensuite des faits
inconnus.\footnote{Vue 62 et vue 63, f. 31r. « le l'esprit méthodique », biffé,
devient « les recherches méthodiques ». « averti » est ajouté dans l'interligne
sans place marquée, et la phrase n'a pas été remise en ordre ; « prévues » est
incertain. Biffé après « conformes à l'expérience » : « en même tems qu'elles
sont explicatives, des faits observés ». Deux mots biffés sont illisibles.}

Aujourd'hui que plusieurs branches de la physique sont entrées dans les
mathématiques, on voit avec admiration les mêmes intégrales, avec des constantes
fournies par des phénomènes de genres différents, représenter des faits entre
lesquels on n'aurait jamais soupçonné d'analogie. Leur ressemblance devient
sensible ; elle est intellectuelle, elle dérive des lois de l'être ; et ce qui
fut le rêve d'une imagination hardie — l'identité des rapports d'ordre et de
proportion dans les existences les plus diverses — se montre aux yeux et à la
pensée avec l'évidence des sciences exactes.\footnote{Vue 63 et haut de la vue
64. Par « les mêmes intégrales », il faut entendre les mêmes équations
différentielles et leurs solutions. L'exemple type, à son époque, est
l'équation de Laplace $\Delta u = 0$, qui régit à la fois le potentiel de la
gravitation, le potentiel électrostatique (Poisson, 1812) et la température
d'équilibre d'un corps (Fourier, 1822) ; l'équation des ondes en est un autre.
L'exemple est de cette lecture ; le feuillet n'en donne aucun.}

\section*{Les lois de l'être dans l'ordre moral et politique : la mécanique des sociétés}

\pagerange{64}{76}
Les lois de l'être ne régissent pas seulement les faits des sciences ; elles
s'appliquent à l'ordre intellectuel. C'est en se rapprochant du type du vrai que
les théories se perfectionnent, que la morale s'épure, que la politique
s'éclaire, que la métaphysique cesse de s'égarer, que les lettres et les arts se
rendent compte de leurs règles. Si, par des progrès qui passent encore toute
espérance raisonnable, la langue du calcul devenait applicable aux questions
morales, politiques, métaphysiques ou de goût, la ressemblance des formules
rendrait évidente celle des objets : leur nature propre serait représentée par
des constantes, et les propositions de chaque sujet par des fonctions dont les
formes reviendraient sans cesse.\footnote{Vue 64 (verso du f. 31 ; le numéro de
page, biffé, ne se lit pas) et vue 65, f. 32r. Le vœu
rejoint ce que Condorcet avait nommé « mathématique sociale » ; la coïncidence
est de cette lecture.}

\subsection*{Les formules, et les causes perturbatrices}

\pagerange{65}{66}
Son exemple : dans plusieurs genres de phénomènes, les formules montrent une
tendance à la régularité, parce que les termes qui expriment l'irrégularité
contiennent le temps de manière à s'évanouir après un temps très court. Le
« théorème » de la courte durée des causes perturbatrices serait alors attesté,
dans les sujets moraux, par les formes du calcul : on verrait combien durent peu
les effets de la fraude, du mensonge et de l'injustice, et que le vrai et le
juste tendent sans cesse à écarter les obstacles ; en politique, on
distinguerait les causes dues à des forces croissantes, qui finissent par
dominer, des causes accidentelles, très fortes un instant et bientôt éteintes ;
dans les sciences de raisonnement, l'erreur tendrait à disparaître ; en matière
de goût, la mode, cause perturbatrice, n'aurait qu'un règne court. « Il est
donc vrai que, quelques divers que soient les sujets, les actions qui troublent
l'ordre naturel tendent à s'anéantir. »\footnote{Vue 65. « Si bien » (le
théorème relatif à la courte durée) est une lecture incertaine. La vue 66, verso
du f. 32, est blanche et porte seulement un nombre biffé.}

Ce théorème n'en est pas un en général, et il faut dire où il est vrai. Dans un
système linéaire \emph{dissipatif} — une oscillation amortie, la conduction de
la chaleur —, la solution se décompose en un régime permanent et un régime
transitoire dont les termes décroissent comme $e^{-kt}$ avec $k > 0$ : là, les
écarts s'éteignent effectivement. Dans un système \emph{conservatif}, rien de
tel : les perturbations mutuelles des planètes, dans la théorie de Lagrange et
de Laplace, sont périodiques et ne s'éteignent pas ; et dans un système
\emph{instable}, elles croissent. La notion moderne est celle de stabilité
asymptotique (Liapounov, 1892).\footnote{Remarque de cette lecture, qui vaut
aussi pour la reprise de la vue 76 : « en toute chose les forces perturbatrices
sont fonctions du tems et \ldots{} la régularité tend à s'établir dans tout
système, de quelque nature qu'il puisse être ». Les sujets moraux et politiques
auxquels elle applique la conclusion n'ont, eux, pas d'équations, et
l'argument y reste une analogie.}

\subsection*{Équilibre stable et équilibre instable}

\pagerange{67}{69}
L'analogie ne se borne pas à un point : la mécanique rationnelle tout entière
offrirait avec la science politique des ressemblances telles que ses théorèmes
deviendraient pour elle des propositions incontestables. Un système de forces
est en équilibre lorsque leur résultante et leur moment résultant sont nuls ;
les forces se compensent, et produisent des réactions dans des directions qui ne
sont pas celles de leur action. De même les forces nées de l'état de société : si
elles se compensent, le repos se maintient de lui-même, et il y a de l'art à
changer la direction où elles agissent en leur opposant des obstacles indirects,
dont le parallélogramme des forces pourrait être l'emblème.\footnote{Vue 67,
f. 33r. Le feuillet dit l'équilibre « vient de ce que l'action des unes est
opposée de direction et égale en puissance à celle des autres » ; la condition
exacte, pour un solide, est la nullité de la somme des forces et de la somme de
leurs moments, et des forces peuvent s'équilibrer sans être deux à deux
opposées.}

Un repos peut être de deux natures. Écarté d'une position d'équilibre
\emph{stable}, un système tend à y revenir et oscille autour d'elle ; les
oscillations s'amortissent s'il y a frottement, et gardent leur amplitude s'il
n'y en a pas. Écarté d'une position d'équilibre \emph{instable}, il s'en éloigne
de plus en plus, et ne retrouve le repos que dans une situation toute
différente, s'il en existe une. Pour un système conservatif, l'équilibre est
stable là où l'énergie potentielle est minimale (théorème de Lagrange–Dirichlet),
et, pour un système pesant, là où le centre de gravité est au plus bas parmi
les positions voisines (principe de Torricelli) ; l'équilibre d'un corps posé
exige en outre que la verticale de son centre de gravité passe par
l'appui.\footnote{Vues 67 à 69 (vue 68, verso du f. 33 ; vue 69, f. 34r). Le
feuillet a « l'équilibre se rétablira au moyen d'oscillations, dont l'amplitude
diminuera à chaque instant », sans condition de frottement ; « L'équilibre
stable a lieu lorsque tous les points du système ont atteint la situation qui
convient à leur tendance naturelle » ; et « L'équilibre d'un système exige que
le centre de gravité soit appuyé. S'il se trouve placé le plus bas possible,
l'équilibre est stable. » Ce dernier énoncé est juste, entendu parmi les
positions voisines. Un mot illisible devant « de repos » (vue 68) ; « en »
(sorte que) incertain ; un astérisque après « qui constitue », sans renvoi.}

Elle retrouve les deux cas dans l'état social : des causes d'agitation produisent
tantôt de légers mouvements qui cessent d'eux-mêmes, tantôt des révolutions qui
ne laissent renaître la paix qu'après de grands changements. La tranquillité
est durable quand chaque individu est dans la situation qui convient à sa
tendance naturelle ; les États gouvernés sans égard aux tendances sociales
restent calmes tant que rien n'agite les esprits, mais la moindre circonstance
les ébranle jusqu'aux fondements, et le mouvement ne cesse que lorsque l'État,
reconstitué sur des bases plus solides, offre à chacun les garanties qu'il
réclamait. Chaque point pesant tend vers le centre de la terre, et la position
qu'il atteint dépend à la fois de sa tendance et de ses liaisons ; chaque
individu tend vers le bien-être, et la première condition est que le bien-être
de chacun nuise le moins possible à celui des autres. Le repos d'un État serait
impossible sans égard à la tendance de l'époque, c'est-à-dire à l'opinion : il
faut lui opposer de puissants obstacles, ou se conformer à ses exigences — et
ce sont là la tranquillité précaire et la tranquillité durable.\footnote{Vue 69.
À la fin du premier paragraphe, d'une grande écriture, un nom souligné lu avec
doute « Gutmann » ; en marge du suivant, une mention encadrée lue avec doute
« Fl. 7 ». Un mot biffé devant « plus solides » est illisible.}

\subsection*{L'impulsion, le centre de gravité et la rotation}

\pagerange{70}{71}
Elle passe aux effets d'une impulsion, et c'est ici que sa mécanique doit être
redressée. Selon le feuillet, une impulsion dirigée par le centre de gravité
meut le système comme si tous ses points étaient réunis en un seul, et toute la
force sert à l'effet voulu ; dirigée à côté, elle se décompose en une portion
qui fait avancer le système et une portion perdue pour ce but, qui le fait
tourner ; et si la première portion était nulle, il ne resterait qu'une rotation,
qui tend à rompre les liaisons du système. Ce qui est vrai est ceci. Le centre
de gravité d'un système se meut comme un point qui porterait toute la masse et
recevrait toutes les forces extérieures, \emph{quelle que soit leur ligne
d'action} (théorème du centre d'inertie) : une impulsion donnée à côté du centre
de gravité lui imprime exactement la même vitesse que si elle passait par lui,
et y ajoute une rotation. Rien n'est retranché au mouvement d'ensemble. Ce qui
est perdu pour lui, c'est de l'énergie : à travail fourni égal, la part dépensée
en rotation manque à la translation. Et une rotation pure, sans avancée du
centre de gravité, ne peut venir d'une force unique : il y faut un couple. Dans
la rotation, enfin, les points les plus éloignés de l'axe vont le plus vite et
demandent la plus grande force pour rester liés, qui est nulle sur
l'axe.\footnote{Vue 70, page 60 barrée, et vue 71 (aucun folio : « 3 » suivi
d'un chiffre ambigu). Le feuillet : « la force seroit décomposée en deux
portions. L'une, celle dont la direction passeroit par le centre de gravité du
système, agiroit comme si elle étoit seule pour faire avancer le système \ldots{}
tandis que l'autre portion de la force motrice, entièrement perdue par rapport à
ce but n'auroit d'autre effet que de faire tourner le système autour de son
centre de gravité » ; puis « si l'impulsion avait été assez maladroite pour que
la première portion de la force motrice fût nulle, le système n'auroit aucun
mouvement progressif » ; et « Les forces tangentielles sont nulles au centre du
système ». Le mot « avait », dans l'interligne, est en lettres droites, d'une
écriture qui ne paraît pas celle du texte.}

L'analogie politique qu'elle en tire garde son sens sous cette forme corrigée,
et c'est à elle qu'on la rapporte : quand l'action du gouvernement suit
l'opinion, la société se meut comme un seul individu qui agirait selon ses
intérêts, et toutes les forces de l'État servent la prospérité ; quand elle la
contrarie en partie, une part de l'effort est dépensée en agitation ; une
administration qui irait en tout contre l'opinion — c'est-à-dire contre
l'intérêt public — agiterait l'État jusqu'à le dissoudre. Les provinces
frontières favoriseraient alors l'État voisin qui voudrait les envahir : en
politique comme en mécanique, les points extrêmes sont les plus agités, et le
désir de séparation serait absurde dans une capitale.

\subsection*{Corps durs, élastiques et mous ; forces vives et levier}

\pagerange{71}{76}
Les sociétés sont faites de trois éléments : des intérêts, des passions, de
l'inertie. Chacun réunit parfois les trois manières d'être, mais l'une domine
d'ordinaire, et ce sont trois caractères, qu'elle compare aux corps durs,
élastiques et mous. Les gens d'intérêt tiennent obstinément leur route et ne
changent de direction que lorsque l'obstacle a épuisé leur force ; les
passionnés prennent au moindre obstacle un parti inattendu, et partent en sens
contraire ; les inertes, amis du repos, souffrent de réelles blessures plutôt que
de réagir. Dans les temps calmes, les intérêts dominent ; ils indiquent
d'eux-mêmes les mesures qui leur conviennent, leur direction est connue, ils
fondent l'opinion publique, et les protéger est aisé. Mais que la paix soit
troublée, les passions accroissent le trouble en agissant dans mille directions,
et l'issue de leur choc est imprévisible. On n'a pas encore fait de statistique
des caractères ; il y a sans doute un grand nombre d'hommes qui agissent toujours
selon leurs intérêts ; le reste se partage entre les êtres irritables, qui
méprisent l'intérêt au prix de leurs passions — dont l'amour-propre est la plus
constante —, et les gens inertes, esclaves de leurs habitudes, sans ambition de
richesse ni de gloire, ni affections vives.\footnote{Vues 71, 72 (verso) et 73
(f. 36r). Le nombre d'hommes d'intérêt est écrit en chiffres et surchargé : « plus
de … sur cent peut-être » ; le chiffre ne se lit pas, et cette lecture n'en donne
aucun. À la fin du paragraphe sur les gens inertes, d'une grande écriture, un
nom souligné lu avec doute « Deraché ».}

Il n'existe, dans la nature, aucun corps parfaitement dur, parfaitement
élastique ou parfaitement mou ; de même, aucun homme n'est intéressé au point de
n'agir jamais par d'autres motifs, les passionnés cèdent parfois à leurs
intérêts, et les inertes peuvent avoir des passions faibles, qui ne sont pas sans
effet. Les trois idéaux mécaniques se disent aujourd'hui par le coefficient de
restitution $e$ d'un choc, rapport des vitesses relatives après et avant : $e =
1$ pour un choc parfaitement élastique, qui conserve l'énergie cinétique, $e = 0$
pour un choc parfaitement mou, où les corps repartent ensemble ; le corps dur
idéal, indéformable, n'a pas de coefficient propre, et la mécanique du
XVIII\textsuperscript{e} siècle le traitait comme un corps mou. Un corps mou
heurté se déplace ; ce qu'il absorbe par sa déformation, c'est une part de
l'énergie cinétique, non la quantité de mouvement, qui se
conserve.\footnote{Vues 73 et 74 (verso du f. 36). Le feuillet définit le corps
parfaitement élastique comme celui « qui ne retienne rien de la direction dans
laquelle on le pousse », et le corps parfaitement mou comme celui « que le choc
ne puisse faire changer de place, et qui absorbe toute la force employée par le
seul changement de forme qu'il subit ». Un mot sous une tache d'encre, après
« subit », est illisible.}

Le trouble d'un État rompt la balance habituelle entre les trois nombres qui
représentent ces caractères : tous reçoivent une impulsion qui les rend
passionnés, la force des masses nouvelles est difficile à évaluer, les
directions sont incertaines, l'agitation se manifeste par des actions
instantanées. En temps de paix, des lumières, l'habitude des affaires et
l'intention de faire le plus de bien avec le moins de mal suffisent à
gouverner ; en temps de crise, il faut se résoudre vite et avoir du courage, qui
ne va pas nécessairement avec l'habileté. S'ajoutent des individus doués par la
nature de grandes forces, que leur position annulait dans le calme, et que le
trouble fait paraître armés d'une énergie inconnue ; sans étude préalable, ils
ne peuvent adopter que les partis violents, où leurs forces trouvent emploi, et
qui dispensent de l'adresse qui leur manque.\footnote{Vue 74 et vue 75, f. 37r.
« Si bien le trouble d'un état rompt la balance » est ainsi au manuscrit.}

« L'égalité est une erreur », écrit-elle, et la mécanique lui fournit encore une
image. Deux masses de même poids peuvent avoir des forces vives très
différentes — la force vive $mv^2$, double de l'énergie cinétique $\tfrac12 mv^2$,
dépend de la vitesse. Et le plus petit poids, placé au bout d'un levier, fait
équilibre à une masse aussi grande qu'on voudra : il suffit que les moments
soient égaux, $P\,a = Q\,b$, ou, ce qui revient au même, que les travaux
virtuels se compensent. De même, les révolutions sont dangereuses et
incertaines parce qu'elles changent d'un coup les rapports entre les forces vives
des différentes classes ; mais, en vertu du même « théorème » que plus haut, ce
qu'elles ont de violent disparaît bientôt.\footnote{Vues 75 et 76 (page 66
barrée). Le feuillet écrit qu'« il ne s'agit que d'établir l'égalité entre les
forces virtuelles » ; ce que la mécanique égale, dans le levier, ce sont les
moments, et, dans le principe de Lagrange, les travaux virtuels (« vitesses
virtuelles »). La conclusion politique est une analogie, et la mécanique ne la
fonde pas : que deux poids égaux puissent avoir des forces vives inégales ne dit
rien de l'égalité entre les personnes. Elle écrit elle-même, trois pages plus
loin, que la tendance à généraliser « a précipité le jugement en avant de
l'expérience ».}

\section*{L'analogie n'est pas une méthode : elle existe par elle-même}

\pagerange{76}{79}
Entre l'ordre physique, l'ordre moral et l'ordre intellectuel, les analogies
sont remarquables, et sans elles le langage figuré n'aurait pu naître. On les a
senties de tout temps, et l'examen des langues a sans doute tout dit à ce sujet ;
il suffit d'observer la justesse de ces passages du sens propre au sens figuré.
La force morale et la force intellectuelle se comportent comme la force
physique : elles se composent et se décomposent, et nos facultés concourent à un
seul acte comme des forces de natures et de directions différentes donnent une
résultante qui représente la vraie force motrice.\footnote{Vue 76 et vue 77,
f. 38r. « Les langues attestent ces analogies », biffé, ouvrait le paragraphe ;
« de ces analogies », écrit au-dessus d'un mot biffé lu avec doute
« desquelles », n'a pas été remis à sa place.}

La justesse de l'expression frappe davantage lorsque, fondée sur un premier
rapport connu, elle vaut encore pour des rapports qu'on n'avait pas en vue. On
appelle monstre un être difforme, et, au moral, un être vicieux, parce que le
vice est une difformité morale. Or, pour elle, l'anatomie a appris depuis que la
monstruosité vient du développement extraordinaire de certains organes, qui
attirent à eux les forces de la vie et privent les autres de leur nourriture ; et
les êtres qui effraient la société par de grands crimes, ou la troublent par des
désordres habituels, ont de même des qualités hors de mesure qui absorbent leur
moralité, et manquent de celles qui, chez des gens moins généreux et moins
courageux, entretiennent l'amour de la justice et de l'ordre. Un seul rapport
bien établi entre deux sujets de genres différents en annonce donc beaucoup
d'autres.\footnote{Vues 77 et 78 (verso du f. 38). Le feuillet : « Lorsque la
langue s'est formée, l'anatomie n'existoit pas. Cette science nous a appris que
la monstruosité avoit pour cause le développement extraordinaire de certains
organes \ldots{} » L'anatomie est bien plus ancienne que ne le dit la phrase
(Galien, Vésale) ; ce qu'elle rapporte est la « loi du balancement des organes »
d'Étienne Geoffroy Saint-Hilaire, énoncée dans sa \emph{Philosophie
anatomique} (1818--1822), au moment où l'étude des monstruosités devenait une
science. La tératologie moderne ne tient pas ce balancement pour la cause des
malformations, qu'elle rapporte à des perturbations du développement,
génétiques ou dues au milieu. Le feuillet ne nomme pas Geoffroy
Saint-Hilaire.}

Cette proposition, énoncée formellement, paraîtrait peut-être hasardée ; chacun
raisonne pourtant comme si elle était indubitable. Elle contient le principe de
l'analogie, qu'« un de nos auteurs » a appelé une méthode d'invention : elle ne
lui refuse pas ce caractère, mais le principe convient aussi aux usages les plus
communs, puisqu'il dérive du sentiment des lois de l'être, partout semblables.
L'habitude de l'étude rend apte à saisir les analogies, et c'est par là qu'elle
facilite les connaissances nouvelles : aux premiers mots d'un sujet inconnu,
l'esprit cherche à quel « module » nouveau appliquer les lois qui conviennent à
tout, puis avance à grands pas, la règle de proportion trouvant mille
applications. Mais il ne faut pas appeler l'analogie une \emph{méthode} :
« L'analogie n'est pas d'invention humaine ; elle existe par elle-même. » Notre
intelligence est faite pour la reconnaître ; elle instruit l'enfant, et l'induit
parfois en erreur, par la même cause qui a produit les systèmes hasardés des
philosophes : partout la tendance à généraliser, née du sentiment intime de
l'unité de l'être, a précipité le jugement en avant de l'expérience, dont il
aurait dû attendre les données.\footnote{Vues 78 et 79 (f. 39r, page 69
barrée). L'auteur qui a appelé l'analogie « une méthode d'invention » n'est pas
nommé, et cette lecture ne le nomme pas. En marge, trois ou quatre mots biffés
et surchargés, illisibles. « méthode » est souligné sur le feuillet.}

\section*{L'état des lettres}

\pagerange{79}{84}
« Il nous reste à présenter quelques considérations sur l'état des lettres aux
diverses époques dont nous avons examiné les opinions systématiques. »\footnote{Vue
79. À la fin du paragraphe précédent, d'une grande écriture, un nom souligné lu
avec doute « Pemeguin », et, au crayon et très pâle, une mention commençant par
« Ch. 5 », dont la suite n'est pas lue. La marque tombe exactement sur ce
changement de sujet.}

Les Anciens, si mal informés des phénomènes naturels, si ignorants des lois qui
les régissent, et pourtant si féconds en généralités qui embrassaient l'univers,
ont atteint la perfection dans tous les genres d'écrire. Il ne faut pas s'en
étonner. Le sentiment du beau est tiré des lois mêmes de la nature intellectuelle
de l'homme ; les observations qu'il demande n'ont besoin ni des instruments des
modernes, ni de cette maturité de raison qui semble avoir remplacé chez eux la
fraîcheur d'imagination, dont la force et la grâce venaient peut-être de son
indépendance. L'art d'émouvoir et celui de plaire ne demandent que la
connaissance des choses humaines. L'esprit s'est d'abord réfléchi sur lui-même :
s'il a pu errer en y cherchant le modèle de l'univers et la cause finale de
toutes choses, il ne pouvait se tromper sur les lois de son propre être. Il était
là, d'emblée, dans la position d'observateur qu'il n'a prise que fort tard à
l'égard des objets extérieurs : les faits intellectuels étaient trop près de lui
pour qu'il ne les vît pas bien. Le génie qui reproduit et transmet des
impressions profondes a donc manifesté sa puissance dès que l'homme a vécu en
société. Le goût est sans doute le fruit d'observations nombreuses et n'a été
fixé que longtemps après les premiers modèles ; mais il n'a pas fallu un temps
immense pour que les remarques fournissent tout ce que l'esprit peut acquérir
dans un genre d'étude exempt, par nature, des causes d'erreur qui ont égaré les
recherches dont l'objet était hors de lui.\footnote{Vue 80 (verso du f. 39 ; le
numéro de page, biffé, n'est pas lu), et haut de la vue 81 (f. 40r, page 71
barrée). Un crochet en marge devant « L'art ». « on sait » et « l'avoient » sont
incertains ; un mot biffé devant « premiers ouvrages » est illisible.}

Les modernes ont voulu imiter la littérature des Anciens et en ont adopté les
fictions, qui ne tenaient plus pour eux à des systèmes ou à des croyances : pleines
de vie entre les mains des inventeurs, elles se décoloraient dans les écrits
d'une nation qui ne les avait pas imaginées, et devenaient énigmatiques et
conventionnelles. Aujourd'hui des formes de style plus conformes au caractère
national s'imposent, et une école nouvelle fait mille efforts pour créer une
littérature qui nous soit propre.\footnote{Vue 81. « à \emph{ses} écrits » :
lecture incertaine. L'« école nouvelle » n'est pas nommée ; ce que la phrase
décrit — le refus de la mythologie antique et la recherche d'une littérature
nationale — est le programme du romantisme français des années 1820, et c'est
de cette lecture que vient le nom.}

Notre époque est aussi marquée par l'invasion des formes mathématiques dans des
ouvrages qui, par nature, ne peuvent atteindre l'exactitude à laquelle ces formes
conviennent : l'emploi maladroit des termes de la certitude produit une sorte de
déception intellectuelle qui choque la raison et le goût. On a reproché aux
géomètres la sécheresse de leur style, sur la foi des seuls éléments ; or les
ouvrages de haute mathématique ont dans le style même un charme inexprimable,
une précision élégante, une extrême finesse, l'art de rendre présentes à la
pensée des idées qui ne sont pas écrites. Tout cela disparaît dans les copies
grotesques qu'on en fait dans le langage commun : on nous présente l'écrin où
nous avons l'habitude de trouver des pierres précieuses, et il contient des
choses de peu de valeur. Le pire est l'emploi des chiffres là où ils n'expriment
aucune valeur réelle : ils usurpent le crédit des connaissances positives et
servent à établir l'erreur en trompant les amis de la vérité. Ceux qui, par
amour des idées exactes, renoncent aux genres faciles pour dévorer la sécheresse
des études élémentaires mériteraient des guides qui aient la conscience du vrai.
Les nations sont aujourd'hui comme les individus qui abordent pour la première
fois des études sérieuses : incapables de juger les ouvrages qu'on leur présente,
elles se dédommagent de leur peine par une confiance aveugle dans les doctrines,
et un mépris des formes qu'avaient adoptées des auteurs moins
solides.\footnote{Vue 81, vue 82 (verso du f. 40) et vue 83 (f. 41r).
« hardiment » (on nous présente) est incertain ; sous « Nous », un petit mot lu
avec doute « Que » ; un mot noirci au-dessus de « a portant », un mot biffé après
« à l'abri » et un autre après « et qui d'ailleurs » sont illisibles ; « ont » et
« qui » sont incertains. « l'enveloppe » du feuillet devient « l'écrin ».}

Le pédantisme était jadis le défaut des savants ; seule, aujourd'hui, la
médiocrité reléguée loin des centres en garde la trace. Mais ce sont des masses
entières qui montrent maintenant une confiance illimitée dans leurs lumières et
un mépris absolu pour ceux qui, fidèles à d'anciens documents, restent étrangers
aux « nouvelles idées » ; la jeunesse surtout, trop instruite à ses propres yeux
pour ne pas dédaigner le ton aimable de plaisanterie qui accompagnait chez nous
l'instruction réelle. La littérature a besoin des idées dominantes pour retenir
l'attention de ses « graves enfans » ; on ne peut plus faire rire qu'aux dépens
des ennemis des innovations, et la raillerie, amère, a perdu sa
grâce.\footnote{Vue 84, verso du f. 41.}

\section*{L'avenir : la politique, les lettres et le parallèle des savoirs}

\pagerange{84}{88}
« Ne nous alarmons pas de ces symptômes : ils ne seront que
transitoires. »\footnote{Vue 84. Avant cet alinéa, deux traits horizontaux à
l'encre dans la marge, et, dans l'interligne, au crayon et d'une autre écriture,
« Ch. 6 » suivi de mots à peine visibles, dont « Avenir de » et un mot
illisible ; dans la marge de gauche, au crayon, « avenir de ».} Le goût du public
pour les idées exactes portera bientôt le talent vers les théories politiques ;
la vérité, trouvant des organes dignes d'elle, paraîtra simple et se reconnaîtra
à ce qu'elle est étrangère à l'emphase des doctrines sentencieuses. La politique
recueillera le petit nombre de vérités qui sont à son usage et adoptera les
formes de sa nature ; comme toute science d'expérience, elle craindra d'énoncer
des théories générales avant d'en avoir vérifié la réalité.\footnote{Vue 85,
f. 42r.}

Une variante, qu'elle a gardée entre crochets, va plus loin : les progrès
immenses dont on fait tant de bruit se réduiront au développement d'idées
contenues dans les ouvrages des prédécesseurs, revêtues des formes nouvelles que
les sciences modernes ont adoptées. Après un temps où une seule partie des
connaissances se distinguait par une méthode sévère, tandis qu'ailleurs les
idées heureuses restaient mêlées aux conjectures hasardées, l'homogénéité qui
caractérisait les travaux des Anciens, dominés en tout par l'imagination, se
retrouvera dans des travaux modernes soumis partout à la marche
méthodique.\footnote{Vue 85. En marge, « Variante », en regard d'un crochet
ouvrant de sa main. Un mot biffé après « le vague et » est illisible.}

Les lettres ont perdu leur éclat ; la poésie, si elle ne touche pas aux
discussions politiques, est dédaignée ; comment le génie y trouverait-il des
inspirations ? Mais un jour plus pur reviendra. L'analogie exige que toutes les
branches du savoir se développent parallèlement : l'attention se portera sur
chacune tour à tour, jusqu'au moment où, leurs progrès devenant comparables,
elles recevront ensemble l'intérêt dû à leur valeur. Ainsi l'élève tourne son
attention d'une étude à l'autre, et semble à chaque fois étranger à ce qu'il
sait déjà, jusqu'au jour où chaque chose reprend à ses yeux son importance, et où
il voit des liaisons là où il ne voyait que des différences. L'esprit humain
touche à cette période : bientôt le tableau des sciences, des lettres et des arts
offrira une symétrie méthodique qui permettra d'embrasser d'un coup d'œil
l'œuvre de l'esprit humain. L'analogie, qui a produit jadis des systèmes
hasardés et des allégories ingénieuses, ne s'arrêtera plus à la surface des
choses : elle en pénétrera la nature, et le type du vrai montrera dans les sujets
les plus divers le caractère commun de toute connaissance
certaine.\footnote{Vue 86, verso du f. 42 (paragraphe ouvert par un crochet), et
vue 87, f. 43r. Une correction retient la pensée : les progrès des branches
devenant « comparables, \emph{ou si j'ose le dire, identiques} » — la seconde
épithète est biffée. Un mot biffé après « leurs progrès » est illisible.}

Les méthodes géométriques, on l'a vu, ont étendu leur empire aux sciences
physiques ; les sciences morales et politiques subiront bientôt la même
transformation. L'opinion publique l'attend, et la devance même par un
enthousiasme irréfléchi pour les doctrines qui la promettent — erreur passagère
et sans danger. Les méthodes existent ; seul l'amour-propre en retarde l'emploi,
les hommes capables craignant de n'être pas estimés de leurs pairs et de ne pas
trouver de juges hors des sciences. L'obstacle ne durera pas, et elle regarde
« dès à présent les sciences morales et politiques comme appartenant au domaine
des idées exactes ». La prévision ne s'est pas réalisée au sens où elle
l'entendait. Le XIX\textsuperscript{e} et le XX\textsuperscript{e} siècle ont
bien mis en nombres et en équations une part de ces sciences — la statistique
sociale, l'économie mathématique, la théorie de la décision —, mais elles n'ont
pas acquis la certitude des sciences exactes, et la politique n'est pas devenue
une science de démonstration.\footnote{Vue 87 (« Laissons les sciences à part »,
biffé, ouvrait le paragraphe ; en marge de droite, de courts traits verticaux,
peut-être le crochet fermant des passages ouverts aux vues 85 et 86) et vue 88,
verso du f. 43. Remarque de cette lecture pour la dernière phrase.}

\section*{Les arts comparés : l'éloquence, la peinture, la musique}

\pagerange{88}{107}
Mais l'analogie montrera toute sa puissance quand l'esprit d'examen comparera la
manière d'agir des lettres et des arts, puis les modèles qu'offre la nature, et y
trouvera partout des copies du modèle du vrai qui appartient à notre être. Ce
modèle, réfléchi autour de nous, cause les impressions que nous recevons des
objets ; le génie produit à son gré des émotions semblables par des moyens aussi
différents que les arts sont séparés, et il doit ce pouvoir au principe de
l'imitation, qui dérive du sentiment d'analogie. Sans un type commun aux objets
qui nous impressionnent, les arts auraient pu copier les objets, les lettres
redire les événements ; ni les uns ni les autres n'auraient su reproduire une
impression semblable à celle des choses réelles par des moyens que le sujet ne
leur fournissait pas.\footnote{Vue 88 et haut de la vue 89 (f. 44r). Page très
corrigée ; l'ordre de lecture est celui qu'impose le sens. Un mot biffé dans
« qu'il sait employer … que constituans » est illisible. À la vue 89, un long
trait en boucle entoure le passage « reproduire \ldots{} sujet », et, en tête de
la page, au crayon et d'une autre écriture, un mot lu avec doute « Double ».}

\subsection*{Le module, l'instant choisi, la gradation}

\pagerange{89}{92}
Les lois de l'être établissent des rapports entre un « module » donné et tout ce
qui tient au sujet dont il est le module, et ce sont ces rapports qui agissent sur
l'imagination. L'âme peut être affectée de la même manière par des sens
différents, parce qu'elle reçoit alors le même ordre d'idées : l'éloquence, la
musique et la peinture procèdent de façon analogue. La peinture ne dispose que
d'un instant, mais elle sait le choisir de manière à rappeler ceux qui précèdent
et à faire pressentir la suite de l'action.\footnote{Vue 89. Le mot « module »
(trois fois ici, puis vues 91, 92, 101 et 110) est clair sur le feuillet, et son
sens n'est pas fixé par la transcription ; dans ces pages il désigne, semble-t-il,
le moyen propre à chaque art et l'échelle qui lui est propre — le son, la
couleur, la parole —, auquel toutes les parties de l'œuvre se rapportent. « Ce
sont » (ces rapports) et, biffé, « pantomime » sont incertains. L'instant choisi
qui rappelle le passé et annonce la suite est, sous le nom de « moment fécond »,
la thèse du \emph{Laocoon} de Lessing (1766) ; le feuillet ne le cite pas.}

Hors des cas où l'auditeur est déjà tout attentif, l'orateur commence avec calme,
et le musicien, de même, ouvre sa pièce par une modulation simple et réserve les
mouvements expressifs pour la suite. La nature a la même gradation : un
crépuscule précède le lever du soleil, un autre en annonce le coucher. Aussi
l'écrivain et l'orateur finissent-ils avec la simplicité du début ; une pensée
éclatante placée à la fin laisserait l'auditeur dans un éblouissement fatigant,
et le goût épuré a proscrit cette manière. Elle affirme que les musiciens, par la
même analogie, renoncent aux accords bruyants qui terminaient naguère les pièces,
et préparent le repos final non seulement par la modulation qui ramène la tonique,
mais par la diminution de l'intensité des sons. C'était une manière, non une
règle : les grandes œuvres de son temps finissent souvent avec éclat, et les
symphonies de Beethoven, pour ne citer qu'elles, s'achèvent presque toutes
\emph{fortissimo}.\footnote{Vue 89 (« aussi », en marge) et vue 90, verso du
f. 44. Le feuillet : « Dans les morceaux modernes, \ldots{} les dernières
mesures préparent le repos absolu ; non pas seulement par la modulation qui
doit amener la tonique, mais encore par diminution d'intensité des sons
employés vers la fin du morceau. » Au-dessus de « qu'elle », quelques mots
biffés et noircis ; « Une pensée » (d'éclat) est sous une tache et incertain ;
un mot biffé après « fatiguant » est illisible.}

Pour percevoir l'identité des rapports qui agissent sur l'imagination dans les
arts qu'on croit étrangers l'un à l'autre, une seule condition est nécessaire :
être familier avec chacun des modules. Les grands effets frappent les moins
instruits, mais la finesse du tact demande de l'exercice ; le goût et l'oreille
se forment. On se croit pourtant incapable d'apprécier l'harmonie, par un
préjugé qui retranche la musique du domaine de l'intelligence. Personne ne veut
paraître insensible aux beautés littéraires, et l'éducation prépare l'homme du
monde à juger, ou du moins à choisir ses autorités ; en musique, l'enfant peu
ému des premiers accords est déclaré non fait pour cet art, l'enseignement est
souvent mauvais, on conclut à l'incapacité, et l'argument de l'inutilité de l'art
achève de le faire abandonner. Si l'on traitait les lettres ainsi, seuls les
êtres d'une sensibilité exquise connaîtraient les chefs-d'œuvre. Les grandes
dispositions sont également rares dans les deux genres ; l'éducation et
l'opinion expliquent seules que tant d'hommes jugent bien de la littérature et
si peu d'une composition musicale. On objecte la disposition naturelle requise
pour juger de la justesse d'un son : mais, même s'il n'était pas sûr que
l'oreille apprend à distinguer des tons qu'elle confondait d'abord, il resterait
une foule de choses, indépendantes du module, qu'un homme intelligent et attentif
pourrait apprécier.\footnote{Vue 90 (deux paragraphes ouverts par un crochet),
vue 91 (f. 45r ; à côté du 45, un petit chiffre à l'encre biffé, peut-être
« 4 ») et vue 92, verso du f. 45. Un mot biffé après « fait » est illisible. À
la vue 92, d'une écriture plus grande, un mot souligné lu avec doute
« Chévenon ». « d'homme capables » est ainsi au manuscrit. Que la
discrimination des hauteurs s'affine par l'exercice est aujourd'hui établi : sa
prudence (« lorsqu'il ne seroit pas certain ») n'était pas nécessaire.}

\subsection*{« La musique est une langue »}

\pagerange{93}{94}
On ne comprend plus ce que l'histoire rapporte de l'influence des différents
modes chez les Anciens ; comment le comprendre, si l'on tient la musique pour
l'art de flatter l'oreille ? Réduite à cela, pourrait-elle retenir une attention
sérieuse, ou produire les effets qu'on raconte ? Ces merveilles cesseront d'étonner
quand on aura montré, en comparant ses moyens à ceux de l'orateur, une vérité que
peu de personnes reconnaissent et qu'elles craignent même d'énoncer : la musique
est une langue, et une langue énergique. Elle emploie les sons, mais les sons ne
la constituent pas. « Elle a ses phrases, ses périodes, ses repos, ses
hardiesses. » Ses beautés flattent l'oreille sans s'y arrêter : elles pénètrent
l'âme et peuvent la gouverner, comme la poésie emploie des sons articulés
agréables, dont on n'aurait qu'une idée incomplète si l'on en oubliait le sens pour
n'en écouter que le nombre.\footnote{Vue 93, f. 46r. Le premier paragraphe de la
page est barré d'une grande croix et inachevé : il parlait de l'importance que
les Anciens attribuaient au « choix des genres », puis à « l'influence
qu'avoit eu le choix des modes », qui semblait attester que « les grecs ont en
effet regardé la musique comme une véritable langue ». La reprise abandonne les
genres pour les seuls modes. La doctrine visée est celle de l'\emph{ethos} des
modes, chez Platon (\emph{République}, III) et Aristote (\emph{Politique},
VIII) ; les « genres » biffés sont les genres diatonique, chromatique et
enharmonique de la théorie grecque. Les noms sont de cette lecture.}

« La musique est toute métaphysique » : ses expressions sont générales, elle n'a
pas de nominatif, elle ne peut exprimer que des sentiments ; mais elle peut mettre
l'âme de l'auditeur dans l'état où l'aurait mise le récit d'une action
particulière. Elle procède comme toutes les langues et fournit des expressions
de tous genres.\footnote{Vue 94, verso du f. 46 (page 84 lue avec doute). « l'a »
(occupée) est incertain. La question de savoir si la musique exprime des
sentiments sera débattue après elle — Hanslick (1854) le niera ; la mention est
de cette lecture.}

\subsection*{La grâce et la gaîté}

\pagerange{94}{96}
Le compositeur doit, comme l'orateur, avoir une idée dominante. Il s'empare de
l'attention par des phrases usitées, que personne ne s'étonne d'entendre, et qui
amènent le sujet sans l'expliquer tout entier. L'écrivain qui veut occuper son
lecteur d'idées gracieuses sans troubler le repos de l'âme sera harmonieux,
rejettera toute expression ambitieuse, sera pur et jamais recherché ; le
compositeur qui vise le même but sera correct et simple, aucune phrase ne
surprendra l'oreille, mais un charme soutenu fera pénétrer dans l'âme le
sentiment de la grâce. L'homme de lettres qui veut égayer saura amener les
contrastes, mais évitera, s'il craint le burlesque, les tours disparates et les
expressions qui conviendraient au tragique. La conversation des personnes
aimables n'a pas besoin d'autres effets pour plaire, et l'on y remarque que la
plupart des choses gracieuses, dites autrement et à peine changées, deviennent
gaies.\footnote{Vue 94 et vue 95 (f. 47r ; la phrase « mais un charme \ldots{}
soutenu » passe d'une page à l'autre). « deviendroient » est écrit en surcharge
et incertain ; au-dessus de « Elles peuvent », le mot « conversations », avec un
renvoi : l'auteur semble avoir voulu « Ces conversations peuvent ».}

En musique, le rapport entre la grâce et la gaîté est sensible : les mêmes modes,
les mêmes coupes de phrases servent aux deux, et le changement de mouvement tient
lieu du changement de manière de dire ; plus de vitesse suffit à rendre gaie une
pièce qui, plus lente, n'eût été que gracieuse. Mais en musique comme en
littérature, la grâce exclut les contrastes qui conviennent à la gaîté, comme la
gaîté, sous peine de burlesque, rejette ceux qui sont réservés aux grandes
émotions.\footnote{Vue 95 (paragraphe ouvert par un crochet) et vue 96, verso du
f. 47. Au-dessus de « en musique comme en littérature », une ligne entièrement
noircie ; un mot biffé devant « été simplement gracieuse » est illisible ; l'ordre
retenu est celui qu'impose le sens.}

\subsection*{La mélancolie et le désespoir}

\pagerange{96}{100}
Les rapports observés entre deux états de bien-être se retrouvent entre des états
de tristesse. La mélancolie et le désespoir peuvent naître d'une même cause, et se
succéder chez la même personne jusqu'à ce que l'impression, s'affaiblissant,
efface l'accent du désespoir et ne laisse qu'une tristesse capable de distraction ;
la mélancolie perd alors son abattement, devient douce, parfois chère à ceux qui
l'éprouvent.\footnote{Vue 96. « également sentis » est incertain, « sentis »
paraissant biffé à son tour. La phrase « jusqu'à ce que \ldots{} de distraction »
ne se construit pas au manuscrit (« peut éloigner ou même fasse disparoître »,
« qu'un susceptible », ce dernier mot incertain) ; la phrase modernisée en donne
le sens probable.}

Ici le texte se rompt. La vue 97 commence par une phrase sans majuscule qui ne se
rattache pas à la dernière phrase, complète, de la vue 96 : « les différens
genres, liés entre eux par une origine commune, inspirent à l'homme de lettres, au
poëte et à l'orateur des compositions qui ont aussi entre elles des rapports
sensibles, et pourtant des caractères divers ».\footnote{Vue 97, f. 48r. La
transcription le constate et ne l'explique pas ; ni la pagination à l'encre (86
et 87 aux vues 96 et 97) ni la foliotation ne signalent de lacune. Il manque
peut-être un début de phrase, peut-être une phrase a-t-elle été déplacée ou
laissée en l'état ; cette lecture ne supplée rien.} Le sens général reste
lisible : genres de sentiments apparentés, compositions apparentées aux
caractères distincts.

La sombre mélancolie exprime les mêmes idées que le désespoir ; dans la
conversation, le passage de l'une à l'autre peut n'être marqué que par le
ralentissement ou la précipitation du discours. Il en va de même dans les
compositions musicales courtes : une complainte d'une sombre tristesse ferait
entendre les accents du désespoir si l'on en pressait le mouvement. Dans les
grandes compositions, au contraire, littéraires ou musicales, les caractères sont
distincts, et une bonne œuvre deviendrait faible ou mauvaise si l'on se contentait
de changer la manière de dire. La raison en est dans notre sensibilité : nous ne
pouvons demeurer longtemps dans la situation violente du désespoir ; elle amène
une fatigue extrême qui, si l'exaltation ne trouble pas entièrement nos facultés,
force l'âme au repos — non le repos du bien-être, mais un abattement proche de la
stupeur, une mélancolie sombre, une tristesse profonde.\footnote{Vue 97
(« complainte » est laissé sans article, « une » ayant été biffé ; paragraphes
ouverts par des crochets) et vue 98, verso du f. 48.}

Toute composition qui veut rendre ces impressions doit donc se plier à notre
sensibilité et faire alterner la peinture de l'extrême tristesse et celle du
désespoir ; ces deux aspects d'un même sentiment se trouvent alors si rapprochés
que la comparaison est inévitable, et il faut marquer entre eux d'autres
différences que celle du mouvement. La nature de ces affections en fournit le
moyen. L'abattement s'exprime lentement, et une sorte de monotonie lui convient :
le poète et l'orateur choisiront des syllabes douces, faciles à prononcer, qui
semblent demander peu d'effort ; le musicien, des notes de même valeur, dont
l'effet approche du silence et nourrit la préoccupation qu'il veut faire sentir.
Cette attention a le double avantage de renforcer l'effet du changement de
manière de dire et de préparer une opposition plus frappante avec les accents du
désespoir. De même, les affections violentes placées à côté d'une sombre douleur
trouveront d'autres nuances que la seule précipitation : des expressions fortes,
des notes dures à l'oreille avertiront l'âme de l'exaltation nouvelle. Le poète,
l'orateur, l'auteur dramatique — écrivain ou musicien — tirent leurs grands effets
de l'art d'amener un mot, une note inattendus : l'âme identifiée à l'action voit
tout à coup un incident qui en redouble l'impression, et se trouble, alarmée d'un
surcroît de malheur dont elle ne peut plus mesurer l'étendue.\footnote{Vue 99
(f. 49r ; en tête, « Georges », à l'encre, dans un paraphe, d'une écriture qui
paraît autre) et vue 100, page 90 barrée (verso du f. 49). La fin de la phrase
sur les notes de même valeur est très remaniée : trois mots d'interligne sont lus
avec doute « mieux », « nourrit », « mieux », le troisième pouvant être
« malheur » ; « nourrit » est retenu ici. « ses » (grands effets) et « alarmée »
sont incertains.}

\subsection*{Le drame lyrique, l'auditeur et l'extrême}

\pagerange{101}{107}
En procédant ainsi, les auteurs suivent l'ordre même des événements qui nous
touchent : tous les jours, un détail nouveau d'un malheur connu en redouble
l'impression jusqu'au désespoir. La ressemblance des arts entre eux et avec les
faits qui nous émeuvent tient à l'identité des rapports, sans laquelle il n'y
aurait pour nous ni sentiment vrai ni idée claire. Chaque art, comme la réalité
qu'il imite, a son module propre, et c'est par là qu'ils diffèrent ; mais le
module adopté, rien n'est plus arbitraire entre les parties de l'action : leur
liaison est nécessaire, et, intervertie, elle détruirait la suite et la
gradation de l'intérêt.\footnote{Vue 101, f. 50r, page 91 soulignée ; au crayon,
d'une autre main, un nom lu avec doute « Pennequin » ; un crochet ouvre la page.
« jetter » (dans le désespoir) et « entr'eux » sont incertains ; des lettres
biffées après « ma » sont illisibles.}

Une composition bien ordonnée, l'âme aime à en suivre le développement ; la
variété prévient la fatigue et soutient l'attention, et plus l'œuvre est forte,
moins elle peut s'arrêter sur une seule impression — la tristesse deviendrait
assoupissante, les accents du désespoir trop répétés ne seraient plus entendus.
La musique, pour qui entend sa langue, est de tous les discours du sentiment celui
qui exige le plus de variété, parce qu'il est le plus pénétrant. La nécessité est
si forte qu'on introduit dans le drame lyrique les incidents les plus
invraisemblables, quand l'action n'en fournit pas, pour donner au compositeur des
sentiments de genres différents à mettre en scène : la partie poétique expose une
action déterminée, la partie musicale saisit l'occasion d'exprimer les sentiments
généraux et abstraits qu'amène le récit.\footnote{Vue 102, page 92, et vue 103,
f. 51r. « stile » est incertain ; biffé, « répand l'a », incertain. Une petite
croix sous la virgule qui suit « compositeur ».}

De là les jugements opposés. L'homme peu sensible à la musique qui va entendre un
chef-d'œuvre lyrique cherche malgré lui une distraction dans l'intrigue ; il
s'étonne du peu de liaison des scènes, ne voit pas l'art des transitions entre des
beautés si énergiques que chacune, prolongée, serait une fatigue horrible, et
sort persuadé d'avoir jugé ce qu'il ne sait pas entendre. L'amateur, lui, est
pleinement satisfait et n'a pas vu les prétendus défauts. Ce jugement est
unanime : âge, sexe, condition, savoir ou ignorance n'y créent aucun désaccord
tant que dure l'impression. L'accord se trouble ensuite ; et, chose inexplicable
à qui ignore la musique, les disputes les plus vives sur les écoles et les
interprètes disparaissent, ou ne se soutiennent plus que par entêtement, en
présence de l'exécution. C'est que la diversité est dans l'opinion, et
l'uniformité dans les facultés de sentir ; l'opinion est une force constante,
l'impression doit être en acte pour valoir, le souvenir la reproduit mal, et elle
s'affaiblit avec l'éloignement de sa cause, laissant l'avantage à l'opinion. Par
la même raison, l'amateur, quelque confiance qu'il ait dans sa raison et son goût
littéraire, n'est pas choqué des invraisemblances : saisi par ce qu'il entend, il
n'a pas le loisir d'en juger l'à-propos.\footnote{Vue 103 (crochet devant « S'il
arrive »), vue 104 (page 94, crochet) et vue 105 (f. 52r). Incertains :
« n'établissent », « pur » (entêtement), « finit par » (dans l'interligne),
« sûrs » ; biffé et incertain, « le ».}

Le compositeur habile, comme les grands orateurs, tient l'âme de ses auditeurs
dans sa main. Quand le trouble qu'il a produit ne peut plus croître, il ne laisse
pas s'affaiblir l'impression qui fait son triomphe : il la remplace par une autre,
et garde ainsi tout son empire. Il mesure les impressions, sait combien de temps
elles peuvent nous posséder, et se garde d'atteindre la limite de nos facultés.
Car ce qui est extrême échappe à l'attention, et la faculté de sentir a ses bornes.
La remarque vaut en tout genre, parce que les lois de l'être sont partout les
mêmes : un malheur extrême n'est pas, dans les premiers moments, senti plus
vivement qu'un malheur moindre ; les instants suivants font reconnaître la
différence. La raison en est que nos jugements ont besoin de connaître leur
objet, et qu'une impression d'une force extrême nous jette dans un trouble qui
rend l'âme incapable de comparer, et donc de juger l'intensité du mal. Mais les
affections qui ont une cause extérieure tendent à s'affaiblir : le chagrin
cuisant semble d'abord devenir plus poignant, parce que l'impression, qui empêchait
d'en mesurer l'étendue, diminue réellement et laisse la réflexion en montrer les
faces ; si la cause a été moins forte, l'effet est connu plus tôt, et
l'affaiblissement apporte un soulagement.\footnote{Vue 106 (page 96 ; en marge de
droite, d'une autre main et d'une encre plus pâle, « Derache ») et vue 107
(f. 53r). Trois mots biffés de l'interligne, au-dessus de « devient », devant
« du mal » et autour de « première a eu une force réelle moins grande », sont
illisibles, ainsi qu'une ligne d'interligne et la ligne suivante, biffées. Dans
la marge de gauche de la vue 107, en regard de « il est dans la nature des
affections » jusqu'à « laisse à la réflexion », un passage d'une quinzaine de
lignes courtes est raturé ligne par ligne et barré d'un trait vertical ; rien
n'en a été lu, et cette lecture ne peut dire s'il complétait ou contredisait le
texte.}

\section*{Conclusion : l'épuisement des lettres et l'époque qui vient}

\pagerange{107}{110}
Les lettres et les arts sont notre ouvrage, et leur but est notre bien-être : leur
première règle est donc de ne pas pousser leur action jusqu'au degré extrême qui
trouble le jugement, ou, s'il est dans leur nature de l'atteindre et que de
grandes beautés en naissent, de nous tirer au plus vite de cette position pénible
en portant notre attention vers une région moins sombre.\footnote{Vue 107 et haut
de la vue 108, page 98.}

On reproche à la littérature son épuisement, et au goût de l'exactitude sa
sécheresse ; il semble que l'imagination perde sa puissance quand la raison
établit son empire. Elle en convient pour le passé. Les temps où des hypothèses
faisaient toute la richesse intellectuelle, où l'homme trouvait en lui-même les
convenances auxquelles il soumettait la nature, étaient très favorables à
l'imagination : les doctrines hasardées ne contrastaient pas avec la science, les
fictions poétiques tiraient leur charme d'une demi-croyance qui leur prêtait une
sorte de réalité, l'histoire et la fable se mêlaient, et ce qui était allégorie
pour les uns était pour les autres le récit de faits merveilleux. Cette
disposition donnait à l'art de bien dire une importance qu'il ne peut garder quand
la condition première est de dire vrai ; la faculté créatrice a disparu avec le
crédit des fictions.\footnote{Vue 108 (« demi » est incertain ; un mot biffé
devant « indécise » est illisible) et vue 109, f. 54r, le dernier folio du
volume, page 99.}

Mais si notre culture attache plus de prix à la solidité qu'à l'éclat, si nous
voulons que la raison domine toutes les productions de l'esprit, si même le goût
des recherches refroidit l'imagination, il ne faut pas désespérer d'une époque où
toutes nos facultés s'uniront dans des productions d'un genre nouveau. Les nations
éprouvent aujourd'hui ce qu'éprouverait un jeune homme passé de la littérature aux
études sérieuses : le charme des premières occupations l'abandonne, une vive
curiosité le remplace, et, l'éducation achevée, chaque chose retrouve à ses yeux
son importance réelle. Et comme l'éducation des sociétés consiste moins à répandre
les connaissances anciennes qu'à en acquérir de nouvelles, comme nous marchons
vers des théories fondées sur des vérités incontestables, nous finirons par
rendre aux branches du savoir une harmonie qu'elles ne durent autrefois qu'à
l'imagination. Tant de vérités de genres différents, groupées autour d'une vérité
première qui est le fait principal du sujet, mettront en pleine lumière
l'identité des rapports entre le module de chaque science et de chaque art et
leurs parties. Les lois de l'être, les conditions du vrai, montrées sous mille
faces à la fois, échaufferont l'imagination ; un enthousiasme nouveau, fondé sur
une base plus solide que celui qui embellit les fictions, inspirera les poètes et
les orateurs. Au lieu de créer l'univers selon les caprices de nos volontés, ils
nous le montreront tel qu'il est ; et si le génie entre dans cette route, il verra
avec admiration que l'art de créer n'a été que celui de copier et de transporter
ailleurs de faibles parties d'un tableau qu'il lui sera donné de savoir peindre
en entier.\footnote{Vue 109 et vue 110, page 100. Biffé à la vue 110 : « saura
nous présenter le charme de cet unité ». Le mot « savoir » est taché d'encre. La
page s'achève sur un point : c'est la dernière page écrite du volume, et le texte
ne se poursuit pas au-delà. Le timbre de la Bibliothèque royale, sur cette vue,
n'est pas transcrit.}

\section*{Ce que les transcriptions laissent en suspens}

\pagerange{7}{110}
Ce qui suit n'est pas la liste des mots douteux — les notes les portent tous —
mais celle des points où la transcription ne permet pas une lecture continue, et
de la façon dont cette lecture les a traités.

\begin{itemize}
  \item \emph{La rupture des vues 96 et 97.} Une phrase complète, puis une phrase
    sans majuscule qui ne s'y rattache pas. Rien ne signale une lacune matérielle ;
    la lecture donne les deux phrases telles quelles et ne supplée rien.
  \item \emph{Les phrases que l'auteur n'a pas remises d'aplomb}, et que la lecture
    donne dans leur sens probable, en le disant en note : la fin de la vue 13 (une
    épithète perdue), la phrase sans verbe de la vue 16, la phrase reconstituée de
    la vue 20 (ordre conjectural et deux mots de l'éditeur), « Le lecteur vient
    ensuite donc cherche » (vue 21), le début de la vue 35 sur Kant (« sa remarque
    expresse »), « et à celles qu'il qu'il doit aux sociétés humaines » (vue 36),
    le « pas » absent de la vue 52, « L'homme, peut-il pas d'autre sujet d'étendue »
    (vue 60), la place d'« averti » (vue 63), la phrase non construite de la vue 96
    et la fin de phrase remaniée de la vue 99.
  \item \emph{Le nombre de la vue 72.} « Plus de … sur cent » : le chiffre,
    surchargé, n'est pas lu, et la lecture n'en avance aucun.
  \item \emph{Le passage marginal de la vue 107}, une quinzaine de lignes raturées
    ligne à ligne et barrées : rien n'en est lu.
  \item \emph{Les additions biffées illisibles}, dont la plus longue est la note
    marginale de six lignes de la vue 23, à l'ouverture du chapitre 2.
  \item \emph{Les chapitres.} Elle n'a titré que les chapitres I et 2. « Ch 4 »,
    « Ch. 5 » et « Ch. 6 » sont au crayon d'une autre main, et aucun « Ch. 3 »
    n'est écrit ; le « II » romain de la vue 33 n'est pas expliqué. La division de
    l'essai en chapitres, au-delà du deuxième, ne peut donc pas être établie à
    partir de ces feuillets, et cette lecture ne l'établit pas.
  \item \emph{Les noms des autres mains} — « Gutmann », « Pennequin », « Derache »
    et leurs variantes, « Varlé », « Cherenoz », « Chévenon », « Georges », et le
    « M. D… » de la vue 33 — sont presque tous lus avec doute ; que ce soient des
    compositeurs est une inférence de la transcription.
  \item \emph{La foliotation} : pas de folio 3 entre les vues 7 et 9, un second
    chiffre ambigu au feuillet de la vue 71 ; et, dans la pagination de sa main,
    ni « 4 » ni la plupart des nombres des vues 61 à 80.
  \item \emph{Le sens de « module »} (vues 10, 78, 89 à 92, 101, 110), que la
    transcription laisse ouvert ; la lecture en propose un en note, comme une
    interprétation.
  \item \emph{Deux renvois sans objet.} La vue 29 promet d'expliquer le paradoxe
    de l'éternité de l'âme « lorsque nous nous occuperons de la liaison établie
    entre la morale et les croyances » ; aucune page du volume ne le fait. Les
    petites croix de renvoi des vues 24 à 40 n'ont pas d'objet sur les pages.
  \item \emph{Le destinataire et la date de la lettre d'envoi} : l'identification
    de Champollion-Figeac et la date de septembre 1833 sont des inférences de cette
    lecture, dites comme telles.
\end{itemize}

\end{document}
